\documentclass[10pt,a4paper]{article}

\usepackage{../preamble-paper}
\addbibresource{../ref.bib}
\AtEveryBibitem{\clearfield{note}}
\numberwithin{equation}{section}

% 目录样式：... 点连接标题与页码
\usepackage{tocloft}
\renewcommand{\cftsecleader}{\cftdotfill{\cftdotsep}}
\renewcommand{\cftsubsecleader}{\cftdotfill{\cftdotsep}}
\renewcommand{\cftdotsep}{4.5}
\renewcommand{\cftsecnumwidth}{2.3em}
% 减少目录内行间距
\setlength{\cftbeforesecskip}{4pt plus 1pt}
\setlength{\cftbeforesubsecskip}{2pt plus 1pt}
\setlength{\cftbeforesubsubsecskip}{1pt plus 0.5pt}
\setlength{\cftbeforeparaskip}{0pt}

% Custom binary operators for transported ring structure
\newcommand{\sqplus}{\mathbin{\boxplus}}
\newcommand{\sqtimes}{\mathbin{\boxtimes}}
\newcommand{\sqminus}{\mathbin{\boxminus}}
\newcommand{\sqdiv}{\mathbin{\boxed{/}}}

\fancyhead[LE,RO]{\thepage}
\fancyhead[RE]{\textbf{Gelfand理论}}
\fancyhead[LO]{\textit{}}

% 覆盖 polyglossia 中文节编号，使用纯数字+点格式（无"第""节""章"）
\titleformat{\section}{\Large\rmfamily\bfseries}{\thesection.}{1em}{}
\titleformat{\subsection}{\large\rmfamily\bfseries}{\thesubsection}{1em}{}
\titleformat{\subsubsection}{\normalsize\rmfamily\bfseries}{\thesubsubsection}{1em}{}
% paragraph 内嵌加粗 + 悬挂缩进 + 间距控制
\newif\ifparablock
\AddToHook{para/begin}{\ifparablock\vspace{1.5ex plus 0.3ex}\parablockfalse\fi}
\titleformat{\paragraph}[runin]{\normalsize\rmfamily\bfseries}{\theparagraph}{1em}{\global\hangindent=2.5em\global\hangafter=1\global\parablocktrue\relax}
\titlespacing{\paragraph}{0pt}{1.5ex plus 0.3ex minus 0.2ex}{0.8em}

\begin{document}

\title{\textbf{Gelfand理论}\\  }
\author{Lao Chen}
\maketitle
\setcounter{tocdepth}{1}
\tableofcontents 

\section{Gelfand对偶}

要讨论\textbf{作用}（action）所构成的数学结构，
也就是对象$X \in \mathcal C$上的自态射集
$\text{End} \ X = \mathcal C(X,X)$，
因恒等态射$1_X: X \to X :: x \mapsto 1_X(x) = x$的存在唯一性，
以及自态射的复合封闭性和结合律，
$\text{End} \ X$自动带有幺半群结构。
Abel群范畴$\textbf{Ab}$引入了交换可逆的加性，
以及对应的模作用和环的概念。
环$R$可以构造$n\times n$维$R$-矩阵环，
其全体记为$R_n^n \simeq \text{End} \ R_n$。
以及集合$X$上的$R$-值函数环
$R^X = \textbf{Set}(X,R)$。
对交换环$R(+,0,\cdot,1)$乘法可逆的要求，
即$R^\times = R \setminus \{ 0 \} \in \textbf{Ab}$。
用（交换）域$k$替代环，
让它们构成$k$-代数，
这样有了非交换的算子代数和交换的函数代数。 

Banach空间$(A,\|\cdot\|)$若是代数，
且范数满足$\| a b\| \leq \|a\| \cdot \|b\|$，
称\textbf{Banach代数}(Banach algebra)。
范数不等式保证了乘法的连续，
即当$a_n \to a,\ b_n \to b$时：
$\|a_n b_n - ab\| 
\leq \|a_nb_n - ab_n\| + \|a b_n - ab\|
\leq \|a_n - a\| \|b_n\| + \|a \|\|b_n - b\| \to 0$。
Banach空间$H$上的线性算子$\mathcal B(H)$是含幺Banach代数。

\begin{framed}
  
$X$是测度空间，
对于实数$p \geq 1$，
可以为其上的函数赋予范数$\|\cdot\|_p$，
并定义$p \geq 1$次可积函数构成的函数空间$L^p(X,\mathbb C)$，
其中的函数$f$满足有限范数：
$$
\|f\|_p = \left( \int_X |f(x)|^p \text dx \right)^\frac{1}{p} < \infty
$$

$p=2$是最特殊的情形，
此时可以拓展到复平方可积函数，
其复范数$\|f\|_2$中的被积项可以通过Hermite内积诱导：
$|f(x)|^2 = f(x) \overline{f(x)}$，
从而有复平方可积空间$L^2(X,\mathbb C)$，
其内积定义为：

\begin{equation}
\langle f, g \rangle = \int_X f(x)\overline{g(x)} \text d x
\label{eq:hilbert-inner-prod}  
\end{equation}

如此构成复Hilbert空间$H = L^2(X,\mathbb C)$。
$\mathcal L(H) = \text{End} \ H$表示$H$上的全体线性算子。
定义算子范数
$ \| A \| = \sup_{\psi \in H} {\| A\psi \|}/{\| \psi \|} $。
若范数$\| A \|$有限称$A$为有界算子．
用$\mathcal B(H)$表示$H$上的全体有界线性算子。
对于可分复Hilbert空间$H$，
$\mathcal B(H)$满足：$\| A B \| \leq \| A \| \| B \|$。
这使得$\mathcal B(H)$构成Banach代数．
矩阵代数是Banach代数，
并复共轭下构成含幺C*-代数．

$n$-维复空间上的有界算子代数就是矩阵代数
$C_n^n = \mathcal B(\mathbb C^n)$．
这为 Lie 群的矩阵表示搭建了框架。
$k^\times$ 描述了 $1$ 维可逆 $k$-线性变换，
而 \eqref{eq:det-wedge} 则对应到了任意维度可逆 $k$-线性变换的充分必要条件：
$\det(T) \in k^\times$。
在矩阵代数 $A = (k_n^n,\circ,I_n)$ 中，
行列式刻画了可逆性：
$A^\times = \{ a \in k_n^n \mid \det a \neq 0 \}$。

\end{framed}

谱理论关注代数元$a \in A$在什么条件下相当于数乘$\lambda \in k$，
使得特征方程$\lambda 1_A = a$成立，
即关于数量的特征函数$k \to k :: \lambda \mapsto p(\lambda) = \lambda 1_A - a = 0$。
在非交换算子代数中它是算子$p(\lambda)$；
在交换函数代数中它是函数$p(\lambda)$。
它们是否退化都由$p(\lambda) = \lambda 1_A - a$的可逆性刻画。
这就要求在含幺代数$(A,1_A)$中讨论。
记 $k^\times = k \setminus \{0\}$ 为域$k$中可逆变换也就是可逆元的全体，
假设$A$是可除代数，
可以通过商代数得到可除代数，
那么 $A^\times = A \setminus \{0\}$ 为可除代数$A$中可逆元的全体，
定义谱集：

\begin{equation}
\sigma(a) = \{ \lambda \in k 
\mid p(\lambda) = \lambda 1_A - a \notin A^\times \}
\label{eq:spec-set}
\end{equation}


为了了解矩阵算子 $a \in A$ 在何种条件下相当于 $\lambda \in k$ 的数乘，
使得方程 $\lambda I_n = a$ 在某种条件下成立，
即要求矩阵算子 $a_\lambda = \lambda I_n - a \notin A^\times$ 不可逆。
在矩阵代数中，
这种可逆性条件是由行列式方程
$p(\lambda) = \det(\lambda I_n - a) = 0$ 刻画的。
行列式代表的楔积有多线性的自然展开，
上式可得到域 $k$ 上的多项式表述 $p(\lambda) \in k[\lambda]$，
称为\textbf{特征多项式}（characteristic polynomial）。
此时体现出复数域 $\mathbb{C}$ 作为代数封闭域的优点，
以后的讨论都在复数域上进行。
令 $A = (\mathbb{C}_n^n,\circ,I_n)$ 为 $n \times n$ 维复矩阵代数，
其特征多项式可以分解为：

$$
p(\lambda) = \det(a_\lambda) = \det(\lambda I_n - a) = \prod_i (\lambda - \lambda_i)
$$

$n \times n$矩阵代数中，
矩阵算子的分量$[p(\lambda)]_i^j = \lambda \delta_i^j - a_i^j$，
在对角构成$n$个一次多项式，
及其对应的作为\textbf{代数集}（algebraic set）的闭点（如果有解）。
借助行列式$\det(p(\lambda))$得到$n$-次特征多项式，
其根是Zariski拓扑意义下这些代数集的并集，
也就是谱集：

\begin{equation}
\sigma(a) = \{ \lambda \in \mathbb{C} \mid a_\lambda = \lambda I_n - a \notin A^\times \} = \{ \lambda_i \}
\end{equation}

$\lambda I_n - a$ 的不可逆性对应了 $p(\lambda) = 0$ 的可解性，
其非退化解 $\lambda \neq 0$ 即为矩阵 $a$ 的特征值。
它提示我们在某种条件下， $a$ 作为算子可以用乘性算子来等效：
$\lambda I_n = a$，
其等效的条件即为特征向量的存在性。

\begin{framed}
  

令 $A = (k_n^n, \circ, I_n)$ 为 $n \times n$ 维 $k$-矩阵代数。
$V = k_n \in \mathbf{Vect}_k$ 是 $n$ 维 $k$-列向量空间，
视为 $A$-左模。
\textbf{外幂算子}（exterior power）$\Lambda^n$ 将其提升到 $n$ 阶外积空间，
同时也把 $T \in \operatorname{End} V$ 提升为 $\Lambda^n(T) \in \operatorname{End} \Lambda^n(V)$：

$$
% https://tikzcd.yichuanshen.de/#N4Igdg9gJgpgziAXAbVABwnAlgFyxMJZABgBpiBdUkANwEMAbAVxiRADUACAXk4GsA+oQC+pdJlz5CKMgCYqtRizZdegkWJAZseAkVnkF9Zq0QgAOuYAydALYAjKHQB6YABTsAlD06X7WAHMAdxgoAJgBLE52EFFxHSl9UnlqY2UzSxsHJ1cPb14-QJCwiKiYuK0JXWkSUgBGIyVTEFjNbUk9FAMG1Ka2VviOmoBmesaTfor26qIAFjHeibMByoTO5ABWQ0X0kD4V6cSULZTFJb3YhVDwhBRQADMAJwhbJDIQHAgkOp3mgBUQNQGHR7DAGAAFKpHECPQIACxwKyeLyQBg+X0QozOu0ydkcLncf08gJAwNBEKhnVJMHuiIqyNeiHen1RvzYuOyBKRz0ZP3RSCxaWaHPxrhJZLBkLW0hh8LpmgZSHm-MQW1JWDAzSgECY9gYrGocJgdCgSDATAYDGoODoWAYbEgmu5KMQADZrRiAOxAjVanV6g0gI0ms0Wq0fW32syO1j0nlIAAcHqQAE42RlzLAcG4ieKQZLKTLYQEEZdhEA
\begin{tikzcd}
V = k_n \arrow[dd, "T"'] \arrow[rr, "\Lambda^n"] &  & \Lambda^n(V) = \bigwedge_i V \arrow[dd, "\Lambda^n(T)"] &                           &    & k \arrow[dd, "\det(T)"'] \\
{} \arrow[rr, Rightarrow]                        &  & {}                                                      & {} \arrow[r, Rightarrow] & {} &                          \\
V = k_n \arrow[rr, "\Lambda^n"']                 &  & \Lambda^n(V) = \bigwedge_i V                            &                           &    & k                        
\end{tikzcd}
$$

根据\eqref{eq:dim-exterior-alg}，
$\dim \Lambda^n(V) = \binom{n}{n} = 1$，
作为 $k$-线性空间有 $\Lambda^n(V) \simeq k$。
在微分几何中，
它对应 Hodge 星算子的维度关系。
$\Lambda^n(T) \in \operatorname{End} \Lambda^n(V)$ 等价于 $\lambda \in \operatorname{End} k$。
$\Lambda^n(T)$ 构成了 $1$ 维的乘性算子，
将其简化为域中的数，
即将行列式定义为：对于 $\forall w \in \Lambda^n(V)$，

\begin{equation}
\det (T) \cdot w = \Lambda^n(T)(w)
\label{eq:det-wedge}
\end{equation}


\end{framed}

有限维线性算子可以用det构造作为Zariski并集的谱集，
有限维行列式依赖于外代数的最高次外幂的$1$-维性质，
而在无限维没有最高次幂，
不存在非零的斜对称$\infty$-线性型。
无限维度没有类似行列式的工具来刻画，
只能用基本的可逆性来定义谱。

延续前面对于矩阵代数的讨论，
若 $A$ 是含幺代数则可以讨论可逆性，
其可逆元全体记为 $A^\times$。
谱理论研究代数元 $a \in A$ 在何种条件下相当于 $\lambda_a \in \mathbb{C}$ 的数量乘法。
矩阵代数通过行列式是否退化来研究，
而在一般的含幺代数中，
任何非零线性泛函 $\varphi \in A^*$ 均满足 $\varphi(1_A)=1$。
$\varphi$ 作用于 $\forall a \in A$ 得到数 $\lambda_a = \varphi(a) \in \mathbb{C}$。
若算子 $a$ 的作用相当于数乘 $\lambda_a 1_A = a$，
则有 $\varphi(a_\lambda) = \varphi(\lambda_a 1_A - a) = \varphi(0_A) = 0$。
若 $a_\lambda = \lambda_a 1_A - a$ 可逆，
则将导出矛盾：
$1 = \varphi(1_A) = \varphi(a_\lambda \cdot a_\lambda^{-1}) = 0$。
故 $a_\lambda = \lambda_a 1_A - a$ 必定不可逆，
此数乘的数落在谱中：
$\lambda_a = \varphi(a) \in \sigma(a)$。

著名的 \textbf{Gleason-Kahane-Zelazko 定理}给出了以下等价条件：
\begin{itemize}
  \item $\varphi$ 是非零且乘性的；
  \item $\varphi(1) = 1$， 且 $\forall a \in A^\times:\ \varphi(a) \neq 0$；
  \item $\forall a \in A:\ \varphi(a) \in \sigma(a)$。
\end{itemize}

即乘性线性泛函的值必定落在相应元素的谱中。
反之，若线性泛函 $\varphi$ 满足 $\varphi(1_A)=1$ 且 $\varphi(a) \in \sigma(a)$ 对所有 $a$ 成立，
则可推出 $\varphi$ 必为乘性泛函。

若$A$是Banach代数则$\sigma(a)$是非空紧集．

复连续函数代数$A = C(X)$是典型的交换代数，
记$Z(f) = \{ x\in X \mid f(x) = 0\}$为$f \in C(X)$的零点集/解集，
$A^\times = \{ f \in C(X) \mid Z(f) = \varnothing \}$．
函数$f\in C(X)$有谱

$$
\begin{aligned}
\sigma(f) 
&= \{ \lambda \in \mathbb C \mid \lambda 1_A - f \notin A^\times \} \\
&= \{ \lambda \in \mathbb C \mid Z(\lambda 1_A - f) \neq \varnothing \}
\end{aligned}
$$

非空的零点集意味着
$\exists \lambda\in\sigma(f),x\in X$，
使得函数$f$在$x$点上相当于幺元按照$\lambda$倍数的缩放．
在函数的值域中
$\forall \lambda \in f(X)$，
$\exists x \in X$
s.t.
$\lambda 1_A(x) = f(x)$，
即
$\lambda 1_A - f \notin A^\times$，
于是$f(X) \subseteq \sigma(f)$．
$\forall \lambda \in \sigma(f)$，
$\exists x \in X$
s.t.
$(\lambda 1_A - f)(x) = 0$，
即 
$\lambda 1_A(x) = \lambda = f(x)$
于是$\sigma(f) \subseteq f(X)$，
得到用零点集描述的谱集：

$$
\sigma(f) = f(X)
$$


\subsection{复代数}

虽然可以继续在域$k$上讨论，
但有许多因素促使我们把接下来的讨论固定到复数域$k = \mathbb C$，
包括但不限于：

\begin{enumerate}
  \item 
  $\mathbb C$是代数封闭域，
  使特征方程对应到特征多项式有根，
  直接得到特征值；
  \item
  $\mathbb C$是$\mathbb R$的二次扩域，
  偶数维的实空间问题可以使用复结构、辛结构、Riemann度量等线性二次型；
  \item
  $\mathbb C$的共轭自动产生了对合，
  从而可以构造C*-代数；
  \item
  $\mathbb C$的解析性质，
  可以利用Cauchy-Riemann公式，
  Liouville定理等，
  得到谱集非空，
  特征多项式有根，
  Gelfand理论等；
  \item
  $\mathbb C$的线性空间、拓扑等结构。
\end{enumerate}

\begin{framed}
  
复函数$f = u + \text i v$的全微分，
可以用两个二元实函数的全微分表示：

\begin{equation}
\text d f 
= \text d u + \text i \text d v
= \left( \dfrac{\partial u}{\partial x} \text d x
+ \dfrac{\partial u}{\partial y}  \text d y \right)
+ \text i \left( \dfrac{\partial v}{\partial x} \text d x
+ \dfrac{\partial v}{\partial y}  \text d y \right)
\end{equation}
\label{eq:combo-total-differencial}

且偏导数有
$
f^\prime_x = \dfrac{\partial u}{\partial x}, \quad
f^\prime_y = \text i \dfrac{\partial v}{\partial y}
$。
正如实分析的基础是实微分的实线性，
复分析的基础类似要求让复微分具有复线性．
在实分析中，
变量 $x$ 和 $y$ 是彼此独立的．
从这里开始，
复分析比二元实分析多了一个非常强的结构，
$\text d x$ 和 $\text d y$ 必须耦合在一起，
要求复函数的全微分可以写为：
$\text d f 
= f^\prime \text d z
= \left(f^\prime_x + \text i f^\prime_y\right)(\text d x + \text i \text d y) 
= \left( \dfrac{\partial u}{\partial x} + \text i \dfrac{\partial v}{\partial y} \right)(\text d x + \text i \text d y) 
= \left( \dfrac{\partial u}{\partial x} \text d x - \frac{\partial v}{\partial y} \text d y \right) 
+ \text i \left( \dfrac{\partial u}{\partial x} \text d y + \frac{\partial v}{\partial y} \text d x \right)$，
对比(\ref{eq:combo-total-differencial})，
复导数$f^\prime$的存在性，
也就是\textbf{全纯(Holomorphic)}的要求，
得到了Cauchy-Riemann方程：

$$
\dfrac{\partial u}{\partial x} = \dfrac{\partial v}{\partial y}, \quad
\dfrac{\partial u}{\partial y} = -\dfrac{\partial v}{\partial x}
$$

实分析中可以轻易找到有界且处处可微但不是常数的函数．
但在复分析中，全纯意味着函数在每个点、每个方向都满足Cauchy-Riemann方程，
全纯性是很强的条件，
它拿走了实际分析中 $x$ 和 $y$ 的独立性，
要求实部和虚部必须以一种相互补偿的方式演化，
$u$ 在 $x$ 方向的增长必须由 $v$ 在 $y$ 方向的增长来平衡．

这种极强的结构约束使得函数无法在保持平滑的同时又被限制在一个范围内波动．
Liouville定理：
$\mathbb C$上的处处全纯有界函数是常数．
有界性约束了函数在全局不能逃逸，
从而在局部也失去了波动的自由，
最终只能退化为常数．

全纯函数在某个点上的导数可以用闭曲线积分表达，
即Cauchy积分公式的导数形式：

$$
f'(z_0) = \frac{1}{2\pi \mathrm{i}} 
\oint \frac{f(z)}{(z - z_0)^2} \, \mathrm{d}z
$$

根据Cauchy估计，
$|f'(z_0)|$极限值是退化的，
意味着$f$在$z_0$是常数，
从而能证明Liouville定理。


\end{framed}

继续关注复含幺代数元$a \in A$的谱集$\sigma(a)$。
若$\sigma(a)$是空集，
那么$p(\lambda) = \lambda 1_A - a$在整个$\mathbb C$上都非退化，
全纯函数$p(\lambda)$的倒数$p(\lambda)^{-1}$在整个$\mathbb C$上
都有定义，
且都是全纯的。
因连续性，
$p(\lambda)^{-1}$在复平面上有界。
根据Liouville定理，
复平面上任何有界\textbf{整函数}（entire function），
也就是$\mathbb C$上处处全纯的函数都是常数。
当$\lambda$向无穷远求极限时，
此极限必须存在且是常数
$\lim_{|\lambda| \to \infty} |p(\lambda)^{-1}| = c$。
若$A$是Banach代数，
其范数收缩性质决定了这个极限的常数必然是$c=0$，
从而要求$\forall \lambda \in \mathbb C:\ p(\lambda)^{-1} = 0$。
这样反证出含幺Banach代数的谱集$\sigma(a)$非空。

\begin{framed}

假设复多项式$p(\lambda)$在$\mathbb C$没有根，
可以定义整函数$p(\lambda)^{-1}$，
$\lim_{|\lambda| \to \infty}|p(\lambda)^{-1}| \to 0$，
类似根据Liouville定理，
$p(\lambda)$必须是常数而无法构成多项式，
因此$p(\lambda)$在$\mathbb C$中有根。
根据因式定理：
若 $\lambda_1$ 是复系数多项式 $p(\lambda)$ 的一个根，
则存在一个唯一的复系数多项式 $q(\lambda)$，其次数为 $n-1$，使得
$p(\lambda) = (\lambda - \lambda_1) q(\lambda)$。
这样可以用归纳法得到所有的根，
从而证明代数基本定理。

在有限维情况，
算子代数的谱集表述为特征多项式的根，
在代数封闭域中谱集非空。

\end{framed}

因$\exists \lambda_a \in \sigma(a)$，
$p(\lambda_a) = \lambda_a 1_A - a \in A^\times$是不可逆的，
在可除Banach代数中只有$0$不可逆，
故$p(\lambda_a) = \lambda_a 1_A - a = 0$，
即$\lambda_a 1_A = a$，
i.e.，
代数元$a$本质上都是复标量$\lambda_a$，
映射
$A \to \mathbb C :: a \to \lambda_a$
是代数同态，
且保持范数等距，
得到Gelfand-Mazur定理。

\subsection{函数代数}

更灵活的应用是，
对于含幺Banach代数$A$，
通过对极大理想$\mathfrak m$做等距同构于复数域的商代数
$A/\mathfrak \simeq \mathbb C$。
令$\text{Max} \ A$为$A$的极大理想构成的集合。
令$A = C(X) = h^\mathbb C X$为某个底空间$X$上的复值函数集，
赋予\textbf{上确界范数}（supremum norm）：
$\|f\|_\infty = \sup_{x\in X} |f(x)|$，
$C(X)$构成交换含幺Banach代数。

$A = C(X)$最基本的极大理想是关于点的退化函数集
$\mathfrak m_x = \{ f \in C(X) \mid f(x) = 0 \}$。
$A$视为线性空间时，
线性泛函构成对偶空间
$\mathbb C^A = A^* = \textbf{Vct}(A,\mathbb C)$。
复代数给复线性空间增强了环乘法结构，
而特征标把这种结构传递到了对偶空间。
根据函数 $a \in A$ 在点 $x \in X$ 上的\textbf{求值映射}（evaluation map）
$A \times X \to \mathbb{C} :: (a,x) \mapsto a(x)$，
诱导：

\begin{equation}
\begin{aligned}
  \Phi: X &\to A^* = \mathbb C^A \\
  x &\mapsto \phi_x \\
\end{aligned}
\qquad
\begin{aligned}
  \phi_x: A &\to \mathbb{C} \\
  a &\mapsto \phi_x(a) = a(x) \\
\end{aligned}
\qquad
\begin{aligned}
  \text{ev}_a: \mathbb C^A &\to \mathbb C \\
  \phi &\mapsto \text{ev}_a(\phi) = \phi(a) \\
\end{aligned}
\label{eq:character-fun-alg}
\end{equation}

由于求值关系的线性性质， $\phi_x$ 显然是线性的：
$\phi_x(\lambda_1 a_1 + \lambda_2 a_2) = \lambda_1 a_1(x) + \lambda_2 a_2(x)$。
此外， $\phi_x(a_1 a_2) = (a_1 a_2)(x) = a_1(x) \cdot a_2(x) = \phi_x(a_1)\phi_x(a_2)$。
在一般性的抽象代数 $A$ 中，
满足这一条件的非退化线性泛函称为\textbf{乘性线性泛函}（multiplicative linear functional）
或者\textbf{特征标}（character）。
特征标的全体构成的集合记为 $\Delta(A)$。


\begin{framed}

Dirac delta函数$\delta$
是一种积分归一化的函数
$\int_\mathbb R \delta(x) \text dx = 1$，
且在$0$以外都退化，
这使得$\delta(0) = +\infty$，
于是它并不能构成普通函数代数的成员，
也不属于可积函数。
略加改造让它适合于任意作为底空间的测度空间$X$，
用筛选性质定义，
使它构成乘性连续线性泛函：

$$
\begin{aligned}
A \times X &\to \mathbb C \\
(a,x_0) &\mapsto \phi_x(a) = \langle \delta_{x_0}, a \rangle 
= \int_X a(x)\delta_{x_0} \text dx = a(x_0) \\
\end{aligned}
$$

从而把求值映射\eqref{eq:max-vanishing-fun}的过程线性化，
这样把点$x \in X$对应到了线性泛函$\phi_x = \langle \delta_x, - \rangle$。

\end{framed}

为了让\eqref{eq:character-fun-alg}的映射
$\text{ev}_a: \mathbb C^A \to \mathbb C$连续，
需要给$A$引入拓扑，
使得对于复开集$\forall V \in \mathcal O_\mathbb C$，
其原像$\text{ev}_a^{-1}(V) = \{ \phi \in \mathbb C^A \mid \phi(a) \in V,\ a \in A, \ V \in \mathcal{O}_{\mathbb{C}} \}$必须是开集，
这样最粗的拓扑称为弱 * 拓扑。
这样构造了Gelfand拓扑$\Delta(A)$。
由于不存在平凡的幂零元，
Gelfand拓扑等价于Zariski拓扑在极大谱上的限制。

作为线性泛函，
特征标 $\phi_x$ 的核就是 $x$ 点退化的函数极大理想：

\begin{equation}
\mathfrak m_x = \ker \phi_x 
= \{ f \in C(X) \mid \phi_x(f) = f(x) = 0 \}
\label{eq:max-vanishing-fun}
\end{equation}

从而诱导了\textbf{剩余域}（residue field）的同构：
$C(X)/\mathfrak{m}_x \simeq \mathbb{C}$。

\paragraph{Gelfand变换的同胚}

\eqref{eq:character-fun-alg}中$\Phi: X \to \Delta(C(X))$的过程是单满的，
拓扑上是同胚。

\paragraph{Gelfand变换的交换含幺C*-代数同构}

当$A$是含幺交换C*-代数时，
$\Delta(A)$在弱 * 拓扑下构成紧Hausdorff空间。
根据$\Delta(A) \times A \to \mathbb C$的自然映射，
类似求值映射得到：
$\hat a : \Delta(A) \to \mathbb C :: \hat a(\phi) = \phi(a)$，
它是拓扑空间上的连续函数$\hat a \in \mathbb C^{\Delta(A)}$，
从而实现了$A \to \mathbb C^{\Delta(A)}$的映射，
称为Gelfand变换。
Gelfand变换对代数结构和对合结构的保持易于证明，
C*-代数的对合和$\|a\|^2 = \|a^*a\|$保证了Gelfand变换是等距变换，
根据Stone-Weistrass定理，
在紧Hausdorff空间上，
Gelfand变换是满射，
于是在含幺C*-代数范畴中，
Gelfand变换构造了等距同构
$A \simeq \mathbb C^{\Delta(A)}$。


当$X$是紧Hausdorff空间时，
所有的特征标都具有前述的构造方式，
特征标和极大理想一一对应，
即$\Delta(C(X)) = \text{Max} \ C(X)$，
并且根据Gelfand-Naimark同胚定理：

$$
X \simeq \Delta(C(X)) = \text{Max} \ C(X)
$$

于是，代数 $A$ 可以被完全视为紧 Hausdorff 谱空间 $\Delta(A)$ 上的连续函数代数来研究。
对于非交换代数，
虽然无法找到经典拓扑空间作为其谱，
但非交换几何允许我们直接把非交换代数本身，
视为某种虚拟的“非交换空间”上的连续函数代数来展开研究。

$$
% https://tikzcd.yichuanshen.de/#N4Igdg9gJgpgziAXAbVABwnAlgFyxMJZABgBpiBdUkANwEMAbAVxiRAB12cYAPHAIwBmwAMIAJOkzgBfENNLpMufIRRkAjFVqMWbABoACALwHOAERgMcdABQibegJSO5CkBmx4CRACylN1PTMrIgg9k6uip4qRGQATFpBuqHmltZ2Ds6R7kpeqsh+CYE6IWGZ2R7K3ih+lMXBbJzcfEKiItIAepxw1gBOFbkxKOrkiSVsA9HVyADMo-XJIHJaMFAA5vBEoIK9EAC2SCMgOBBIcQulABZd7Ht0OJf8-AYiINQMdPyWAAqD1SAMGCCHDZHb7JB+Y6nRBzbQNFLsCxWOimdjYPYwACOqOaOGAAFk6DxZO9Pj8-qoAUCQfJtrsDjDqCdDrSQGCGecoRDWeykAA2JnQgDsF0aaKwGMxy2kQA
\begin{tikzcd}
\textbf{CHaus}                               & {} \arrow[rr, "\simeq"] &  & {} & \textbf{CC}^\star                             \\
X = \Delta(C(X)) \arrow[rrrr, "h^\mathbb C"] &                         &  &    & C(X) \arrow[d]                                \\
\Delta(C(X)) \arrow[u]                       &                         &  &    & C(X) \arrow[llll, "\Delta \simeq \text{Max}"]
\end{tikzcd}
$$

\subsection{算子代数}

交换代数，
例如上面讨论的函数代数，
以及代数几何中讨论的整数代数、多项式代数等，
都有丰富的内部结构，
体现为不同形式的理想、素理想、极大理想，
和对应的商代数。
$C(X)$的极大理想或者特征标如此之丰富，
在其上可以构造完整的底空间$X$的拓扑信息。
由于极大理想足够多，
能够按照$T_2$分离任意元素。
交换代数中的理想，
在包含关系中构成偏序，
从而可以构造拓扑，
如Zariski拓扑或Gelfand拓扑，
进而在拓扑上构造层。

令人惊奇的是，
哪怕是最简单的非交换代数，
例如$n > 1$的$X = \mathbb C$上的矩阵代数$\mathbb C_n^n \simeq \text{End} \ \mathbb C_n$，
它都没有非平凡的特征标、极大理想，
当然也就没有拓扑。
在非交换代数中，
没有理想就没法分离点，
没有可以做$T_2$的Hausdorff空间。
非交换代数的理想需要区分左右双边，
可分离性所需要的极大双边理想是稀缺的。
从算子代数角度看，
还存在各种非平凡的投影、幂等元，
它们在代数内部构成了格结构，
产生了复杂的内部非对易关系。

交换代数以理想或者特征标作为点来构造拓扑空间的模式，
在非交换代数彻底失效了。
我们需要引入范畴的思维方式，
用动态的箭头来替代静态的点。
接下来我们先研究量子理论的空间Hilbert空间的性质，
在其中提取一些内部结构，
以便之后构造Heyting代数、
层范畴、
topos。


\section{Hilbert空间}

延续\eqref{eq:hilbert-inner-prod}的记号，
Hilbert空间$H$的自伴算子和有界自伴算子的全体分别记为
$\mathcal L(H)^\dagger$和$\mathcal B(H)^\dagger$，
构成实线性空间．
可观测量对应$\mathcal B(H)^\dagger$中的自伴算子．
算子$P\in \mathcal L(H)^\dagger$若满足
$P^2 = P$称为正交投影算子，
简称投影算子，
非正交的投影算子不是自伴算子，
不在讨论范围．
投影算子的全体记$P(H) \subseteq \mathcal L(H)^\dagger$．
投影算子$P,Q\in P(H)$可等价地定义偏序$\preceq$：

$$
P \preceq Q 
\quad \Longleftrightarrow \quad 
\text{Im} \ P \subseteq \text{Im} \ Q 
\quad \Longleftrightarrow \quad 
PQ = P = QP
$$
 
有偏序范畴$(P(H),\preceq)$，
它是完备格，
最大元$1_H$对应全空间$H$，
最小元$0_H$对应平凡空间$\{0\}$．
投影算子$P\in P(H)$作用于$H$，
使得像闭子空间$PH$产生了包含偏序．
$H$的闭子空间全体构成格$(\text{Sub} \ H,\subseteq)$：

$$
\begin{aligned}
  (P(H),\preceq) &\to (\text{Sub} \ H,\subseteq) \\
  1_H &\mapsto 1_H H = H \\
  0_H &\mapsto 0_H H = \{0\} \\
  (P \preceq Q ) &\mapsto (PH \subseteq QH)
\end{aligned}
$$

实际上在正交模格范畴中
$(P(H),\preceq,1_H,0_H) \simeq (\text{Sub} \ H,\subseteq,H,\{0\})$．
它们不是分配格，
从而不是Bool代数．
$\mathcal B(H)$中的交换von Neumann子代数全体记为$\mathcal V(H)$．
对于$V\in\mathcal V(H)$，
其交换性决定了$P(V) = P(H) \cap V$是$P(H)$的Bool子格．
在context $V$内部逻辑是经典的．

Hermite矩阵$A \in \mathbb C_n^n$相当于矩阵代数中的自伴算子，
其对角化相当于谱分解：

\begin{equation}
A = U \Lambda U^\dagger = \text{diag}(\{\lambda_i\})
= \sum_i \lambda_i P_i
\label{eq:spectral_decomp-finite}
\end{equation}

其中$U$是酉矩阵，
$\lambda_i \in \mathbb R$是实特征值．
$E_i^i$是第$i$个对角元为$1$的矩阵，
$P_i = U E_i^i U^\dagger \in P(\mathbb C_n^n)$是投影算子，
投影到第$i$个特征向量．
在一般的复Hilbert空间$H \in \textbf{Hilb}$上，
把Borel σ-代数$B(\sigma(A))$中的每个可测集$E$映射为$H$上的一个正交投影算子$\mu(E)$：
$\mu: B(\sigma(A)) \to P(H) :: E \mapsto \mu(E)$。
若它是正交模格同态，
则构成了投影值测度。
若视这些集合$B(\sigma(A))$，$\text{Sub} \ H$，$P(H)$为范畴，
则$\mu$是一个函子，
它在有限维中有特征值分解\eqref{eq:spectral_decomp-finite}。
谱定理，
对有界自伴算子$A\in \mathcal B(H)^\dagger$
存在唯一的投影值测度$\mu^A$，
定义在谱$\sigma(A)$上的Borel σ-代数上，
使得：

\begin{equation}
A = \int_{\sigma(A)} \lambda \text d\mu^A(\lambda)
\label{eq:spectral_decomp}
\end{equation}


\subsection{闭子空间}

复Hilbert空间$H \in \textbf{Hilb}$
的内积结构自动赋予了范数拓扑，
定义其子对象时要考虑拓扑性质。
若$H$的子空间在范数拓扑下是闭集，
称为闭子空间．
有限维中所有子空间都是闭子空间，
无限维中闭子空间有特别的性质，
它包含所有的极限点，
易于做分析．

$H$的所有闭子空间构成的集合记为 $\text{Sub} \ H$，
这里的记号暗示了它是一个子对象格：
以包含关系为偏序，
构成正交模格$(\text{Sub} \ H,\subseteq,H,0)$．

\begin{framed}

它不满足分配律，
不是Heyting代数，
更不是Boolean代数．
由于没有子对象分类器，
它不构成topos．
实际上$\text{Sub} \ H$对应着量子逻辑，
而非直觉主义逻辑．

由于$\text{Sub} \ H$十分接近一个topos，
人们试图用各种手段将其引入到topos理论中．

为交换von Neumann代数范畴构造Bohr topos，
其子对象格恢复为$\text{Sub} \ H$．

由于它是正交模格，
可以嵌入到一个Heyting代数中．

我们可以在$\text{Sub} \ H$上构造一个Grothendieck topos，
使得$\text{Sub} \ H$成为这个tope的子对象格．
  
\end{framed}

复Hilbert空间$H \in \textbf{Hilb}$上按照包含可以定义两个偏序范畴：
闭子空间正交模格$\text{Sub} \ H$和投影算子集合 $\text{Proj} \ H$，
它们作为范畴同构：

$$
\text{Sub} \ H \simeq \text{Proj} \ H
$$

对有界自伴算子$A\in \mathcal B(H)^\dagger$，
其谱集$\sigma(A) \subset \mathbb R$上的Borel σ-代数$B(\sigma(A))$，
作为谱集的子集族，
在包含关系下构成Boolean代数．
作为正交模格，
它与$H$的闭子空间范畴$\text{Sub} \ H$之间存在自然同构：

$$
B(\sigma(A)) \simeq \text{Sub} \ H
$$

这样便建立了正交模格的同构：

$$
\begin{aligned}
  P:B(\sigma(A)) &\to \text{Proj} \ H \\
  E &\mapsto P_E
\end{aligned}
$$


\subsection{von Neumann代数}

复Hilbert空间上有界线性算子代数$\mathcal B(H)$，
其含幺子代数$V$称为von Neumann代数，
若它对伴随封闭，
对弱算子拓扑封闭，
即双交换子等于自身。

对于一组有界算子$S \subset \mathcal B(H)$，
其\textbf{交换子}（commutant）为：

$$
S^\prime = \{ T \in \mathcal B(H) \mid \forall A \in S:\ TA = AT \}
$$

如果$M \subset \mathcal B(H)$是子代数，
且$M = M^{\prime\prime}$，
称为\textbf{von Neumann代数}（von Neumann algebra），
它在弱算子拓扑下是闭的。
\cite{MurrayVonNeumann1936RO1}

\subsection{经典力学与量子力学}

以下用实矩阵算子表示的复数，
体现$\dim_\mathbb R \mathbb C = 2$，
因$\{I,J\}$是$\mathbb C$的正交自然基：

$$
z = x + \text i y 
= \begin{bmatrix}
  x & - y \\ y & x
\end{bmatrix}
= x I + y J
= x\begin{bmatrix}
  1 & 0 \\ 0 & 1
\end{bmatrix}
+ y\begin{bmatrix}
  0 & - 1 \\ 1 & 0
\end{bmatrix}
$$

这里的$J$相当于虚数单位$\text i$。
扩展到$n\times n$维复矩阵$[z_i^j] = [x_i^j] + \text i[y_i^j]$：

$$
[z_i^j]
=
\begin{bmatrix}
  [x_i^j] & [-y_i^j] \\
  [y_i^j] & [x_i^j]
\end{bmatrix}
= [x_i^j] I_{2n} + [y_i^j] J_{2n}
=
[x_i^j] \begin{bmatrix}
  I_n & 0 \\
  0 & I_n
\end{bmatrix}
+
[y_i^j] \begin{bmatrix}
  0 & -I_n \\
  I_n & 0
\end{bmatrix}
$$

这里出现了偶数维的单位矩阵$I_{2n}$和辛矩阵$J_{2n}$。
一般考虑实域，
$(Z,\Omega)$称为辛向量空间，
若$\Omega$是$Z$非退化斜对称的双线性型，
即辛形式。
在一对对偶空间$Z = V^* \times V$上构造
\textbf{正则辛形式}（canonical symplectic form）：

\begin{equation}
\begin{aligned}
  \Omega: Z = V^* \times V &\to \mathbb R \\
  ((f,u),(g,v)) &\mapsto 
  \Omega((f,u),(g,v))
  = g(u) - f(v)
  = \langle g,u \rangle - \langle f,v \rangle
\end{aligned}
\label{eq:canonical-symp-form}  
\end{equation}

在有限维$V \simeq \mathbb R_n$时，
$\Omega: \mathbb R^n \times \mathbb R_n \to \mathbb R$
有坐标表示：

$$
((f,u),(g,v)) \mapsto 
  \Omega((f,u),(g,v))
  = f_2 v_1 - f_1 v_2
  = 
  \begin{bmatrix} f & g \end{bmatrix}
  J_{2n}
  \begin{bmatrix} u \\ v \end{bmatrix}
$$

辛形式作为$(1,1)$-型张量，
其坐标表示就是辛矩阵$J_{2n}$。
这类张量易于转为内积，
从而在内积空间上可以定义辛形式：
$\Omega((f,u),(g,v)) = \langle f,u \rangle - \langle g,v \rangle$。
复数以共轭乘法构成内积空间：

$$
\begin{aligned}
\langle z,w \rangle &= \overline z w = (xu + yv) + (-xv + yu) \text i \\
&=\begin{bmatrix} 
  xu + yv & xv - yu \\ 
  -xv + yu & xu + yv
\end{bmatrix} \\
& = (xu + yv) I + (-xv + yu) J
\end{aligned}
$$

这里的实部是实内积，虚部是辛形式．
Hermite内积实际上给出了两个复数的夹角，
保持了旋转不变性．
推广到复矩阵代数的Hermite内积
（Einstein求和约定）：

$$
\langle [z_i],[w_i] \rangle 
= z_i \overline w^j 
= (x_ju^j + y_jv^j) + \text i(u_jy^j + v_jx^j)
= \Re \langle z,w \rangle - \Im\langle z,w \rangle
$$

实部是实内积，
虚部是辛内积。
这样把最常见的二次型张量：
Riemann度量，
复结构，
辛结构统一起来了，
推广到流形的丛就是Kähler流形。
复矩阵代数再推广到Hilbert空间，
得到量子力学的辛形式
\cite{MarsdenRatiu1994}：
$\Omega(\psi_1,\psi_2) = - 2 \hbar h \Im\langle \psi_1,\psi_2 \rangle$。


\subsection{Hermite内积}


$k$-线性空间$H \in \textbf{Vct}_k$有自然配对
$\textbf{Vct}(H^* \times H,k)$，
它是$k$-双线性型．
Hilbert空间$H$是完备的内积空间，
内积是$\textbf{Vct}(H \times H,k)$形式的是$k$-双线性型．
Riesz引理将这两种双线性型对应起来．

限制在复Hilbert空间范畴$\textbf{Hilb}$上讨论，
底空间$X$上的平方可积函数全体记：

$$
L^2(X) = \left\{ \psi :\ \int_X |\psi(x)|^2 \text dx < +\infty  \right\}
$$

因其平方可积，
按逐点Hermite内积的方式定义内积：


\section{Topos}

回到非交换的算子代数$A$，
它是我们研究的非交换的全局空间。
关注其中某个交换子代数$C \subset A$，
它对应了一个经典的\textbf{上下文}（context），
提供了局部信息，
且作为交换代数有自己的Gelfand谱$\Sigma_C$，
可以做经典的点和拓扑。
$A$的交换子代数的全体记为$\mathcal C(A)$，
这一记号体现了它是一个范畴，
其对象为交换子代数，
态射为包含$C \hookrightarrow C^\prime$，
于是在底范畴$\mathcal C(A)$上可以构造层
$\underline\Sigma\in\textbf{Sh}(\mathcal C(A))$：

$$
% https://tikzcd.yichuanshen.de/#N4Igdg9gJgpgziAXAbVABwnAlgFyxMJZABgBpiBdUkANwEMAbAVxiRAB12BbOnACwDGjAAQBhABQBBAJQA9TjhgAPHMAhoAviA2l0mXPkIoyARiq1GLNqO26QGbHgJEyAZnP1mrRCFHz2aABOWFysOnqOhkQATOQelt4c7IoqAEYAZsAAKupa4fb6TkbIsWbUnlY+CsqqALJ0ShrCnGK2EQbOKLHu5Qls1SrA9Y3N7GL+QSFhdg4dxWTR8V5sbQWRnSWki73LPtrmMFAA5vBEoOmBEFxIJtQ4EEixIDh0WAxsfBAQANarF1c3O4PRAAFny-2uiDIz2Brh2lSSTDAsECDCwYBgnAAylgjjw-pdIU97kgAKzgwlkoFIEHwxIDVSBeBaagMOipGAMAAKhSiPmCRz4OAJAMQADZqYgAOys9GJKAQJipBisah8GB0KBIMBMBgMO6vd4+SAY-YaIA
\begin{tikzcd}
\mathcal C(A)^\text{op} \arrow[rr, "\underline\Sigma"] &  & \textbf{Top}                                    \\
C \arrow[dd, hook] \arrow[rr]                          &  & \text{Max} \ C                                  \\
{} \arrow[rr, Rightarrow]                              &  & {}                                              \\
C^\prime \arrow[rr]                                    &  & \text{Max} \ C^\prime \arrow[uu, "\text{res}"']
\end{tikzcd}
$$

\begin{framed}

$n > 1$的复矩阵代数$\mathbb C_n^n \simeq \text{End} \ \mathbb C_n$，
它的交换子代数是可对角化的矩阵族。

\end{framed}

层范畴$\textbf{Sh}(\mathcal C(A))$是一个topos，
点是函子，
逻辑是Heyting代数，
交换子代数的开集格之并。
Topos的子对象是层$\underline\Sigma$的局部截面，
而单个$C$只是截面定义域的一部分。

层的\textbf{全局截面}（global section）$\gamma$可以把大的上下文，
以相容的方式限制到小的上下文。
如果存在全局截面，
所有上下文共享一个相容的点集，
则可以为所有的经典视角上下文一致性地指派一个态，
这正是一个非上下文的隐变量模型。

Kochen–Specker定理说明，
这种局部截面无法拼接成全局截面
~\cite{KochenSpecker1967}：
令$H$为Hilbert空间，
当$\dim H \geq 3$时，
不可能为每个可观测量（自伴算子）同时赋予一个确定的测量值，
使得这些值在所有交换的可观测量子集上，
满足函数依赖关系（即 $f(A)$ 的值等于 $f$ 作用于 $A$ 的值）。
换句话说：
你可以在每个交换子代数（即每个“局部坐标系”）上定义一个经典赋值（截面），
但你无法把这些局部赋值粘合起来，
形成一个全局赋值（全局截面）。


\subsection{Topos}

预层$F: \mathcal C^\text{op} \to \textbf{Set}$是函子，
预层构成的范畴$\textbf{Set}^{\mathcal C^\text{op}}$是函子范畴。
构造到预层范畴的形状为$\mathcal D$的\textbf{图}（diagram），
即从索引范畴到目标范畴的函子
$J: \mathcal D \to \textbf{Set}^{\mathcal C^\text{op}} $，
使得$\forall d \in D$得到反变函子
$J(d): \mathcal C^\text{op} \to \textbf{Set}$。
其极限定义为：
$$
(\lim_{d \in \mathcal{D}} J(d))(C) := \lim_{d \in \mathcal{D}} (J(d)(C)).
$$
由于 \(\mathbf{Set}\) 是完备的，右侧极限存在。
且该构造自然满足函子性，
故$\textbf{Set}^{\mathcal C^\text{op}}$的极限存在。
有终对象
$\mathbf{1}(C) = \{*\}$（恒为单点集的预层），
以及拉回：逐点取集合的拉回。

预层$\forall X \in \textbf{Set}^{\mathcal C^\text{op}}$，
定义幂对象：
$$
\mathcal{P}(X)(C) := \{ \text{子预层 } S 
\subseteq \text{hom}_{\mathcal{C}}(-, C) \times X \}.
$$
即，它是对象 \(C\) 所代表的预层与 \(X\) 的乘积的所有子预层构成的集合。

对于态射 \(f: D \to C\)，限制映射为 \(\mathcal{P}(X)(f): \mathcal{P}(X)(C) \to \mathcal{P}(X)(D)\)，
将 \(S\) 拉回至 \(D\)：\(S \mapsto S \times_{\text{hom}(-, C)} \text{hom}(-, D)\)。

\textbf{关键验证：} 我们需要建立自然同构。
对于任意预层 \(Y\)，一个自然变换 \(\alpha: Y \to \mathcal{P}(X)\) 等价于在 \(\mathcal{C}\) 的每个对象 \(C\) 处选择 \(Y(C)\) 的一个子预层 \(S_C \subseteq \text{hom}(-, C) \times X\)。这又等价于给出一个子预层 \(S \subseteq Y \times X\)。因此，
\[
\text{hom}(Y, \mathcal{P}(X)) \cong \text{Sub}(Y \times X).
\]
故 \(\mathcal{P}(X)\) 是 \(X\) 的幂对象。


初等 Topos 要求对于任意对象 \(Y, Z\)，
存在指数对象 \(Z^Y\) 使得 \(\text{hom}(X \times Y, Z) \cong \text{hom}(X, Z^Y)\)。

预层范畴具有指数对象。

对于预层 \(Y, Z \in \mathbf{Set}^{\mathcal{C}^{op}}\)，定义指数预层 \(Z^Y\) 为：
\[
Z^Y(C) := \text{hom}_{\mathbf{Set}^{\mathcal{C}^{op}}}(\text{hom}(-, C) \times Y, Z).
\]
这本质上是将 \(Z\) 提升为 \(Y\) 的“内部 Hom”。可以验证，存在求值态射 \(\mathrm{ev}: Z^Y \times Y \to Z\)，并且该构造给出了所需的伴随关系。


这是证明中最关键的一步。

预层范畴存在子对象分类器（Subobject Classifier）。

定义预层 \(\Omega \in \mathbf{Set}^{\mathcal{C}^{op}}\) 如下：
\[
\Omega(C) := \{ S \subseteq \text{hom}_{\mathcal{C}}(-, C) \mid S \text{ 是一个筛 (Sieve)} \}.
\]
即，\(\Omega(C)\) 是对象 \(C\) 上的所有筛构成的集合。

对于态射 \(f: D \to C\)，限制映射为：
\[
\Omega(f): \Omega(C) \to \Omega(D), \quad S \mapsto f^*(S) = \{ g: E \to D \mid f \circ g \in S \}.
\]

定义真值态射（Truth Morphism）\(\top: \mathbf{1} \to \Omega\) 为：
\[
\top_C: \{*\} \to \Omega(C), \quad * \mapsto \max S_C = \text{hom}_{\mathcal{C}}(-, C).
\]
即，它将每个对象映射到该对象上的最大筛（整个 Hom 集）。

现在，对任意单态射 \(i: X \hookrightarrow Y\)，我们定义其特征态射 \(\chi_i: Y \to \Omega\) 如下：
对于任意 \(C \in \mathcal{C}\) 和 \(y \in Y(C)\)，
\[
\chi_i(C)(y) := \{ f: D \to C \mid Y(f)(y) \in X(D) \} \in \Omega(C).
\]
即，它是所有使得 \(y\) 沿 \(f\) 拉回后落到子对象 \(X\) 中的态射 \(f\) 构成的集合。这显然是一个筛。

由此可得，对任意 \(Y\) 的子对象 \(X \hookrightarrow Y\)，存在唯一的态射 \(\chi\) 使得下图是拉回：
\[
\begin{tikzcd}
X \arrow[r] \arrow[d] & \mathbf{1} \arrow[d, "\top"] \\
Y \arrow[r, "\chi"] & \Omega
\end{tikzcd}
\]
其中 \(X\) 是 \(Y\) 通过 \(\chi\) 拉回 \(\top\) 得到的。因此，\(\Omega\) 是子对象分类器。

它是初等topos，
当$\mathcal C \in \textbf{cat}$是小范畴时，
它是Grothendieck topos。

\subsection{拓扑}

集合$X \in \textbf{Set}$给出了幂集$PX$，
$PX$的子集为$X$的子集族．
$\{\varnothing,X\}$和$PX$都是子集族，
构成平凡拓扑和离散拓扑．
其它的拓扑空间$\mathcal O_X$介于这二者之间：

$$
\{\varnothing,X\} \subset \mathcal O_X \subset PX
$$

回看拓扑公理对有限交的要求，
若要求无限交，
开集的空间就会被挤压坍塌为单点，
如果单点都是开集，
那么拓扑就成为离散拓扑，
失去了拓扑的信息．

\subsection{Locale}

拓扑空间的开集族$\mathcal O_X$带有偏序结构，
且根据拓扑公理在有限交和任意并上封闭，
并满足无限分配律

$$
U \wedge \bigvee_i V_i = \bigvee_i (U \wedge V_i)
$$

满足这种有限交和任意并封闭，
以及无限分配律的格，
称为frame．
Frame的范畴$\textbf{Frm}$中，
态射为保持有限交任意并的偏序集映射．

$$
% https://tikzcd.yichuanshen.de/#N4Igdg9gJgpgziAXAbVABwnAlgFyxMJZAJgBoAGAXVJADcBDAGwFcYkQBlEAX1PU1z5CKAGwVqdJq3YAVHnxAZseAkTLEJDFm0QgAOnoC29HAAsAxkwAEAeQD6XXv2VCiYjTS3TdB42cuMtnZyTooCKsLI5OKeUjr6ejgwAB44AEYAZsAyEGjc8s6CqijRHpLa7AZJqZnAAGIATob53BIwUADm8ESgGQ0QhkjRIDgQSACMsRW6GQUgfQNIAMw0o0hk5d7zAHrAALTjLZTcQA
\begin{tikzcd}
\textbf{Top} &  & S \arrow[rrrr, "f"] &  &  &  & T                                   \\
             &  &                     &  &  &  &                                     \\
\textbf{Frm} &  & \mathcal O_S        &  &  &  & \mathcal O_T \arrow[llll, "f^{-1}"]
\end{tikzcd}
$$

\subsection{预层}

范畴$\mathcal C$上的预层$\textbf{Set}^{\mathcal C^\text{op}}$是取值在集合范畴上的反变函子范畴．
从最基本的拓扑空间范畴$\textbf{Top}$构造几何结构，
以拓扑空间$X\in\textbf{Top}$中的开集为对象，
以开集的包含关系$U \stackrel{\subseteq}{\longrightarrow} V$为态射，
构成范畴$\mathcal C = \mathcal O_X$，
称为开集范畴，
其预层范畴为$\textbf{PSh}_\textbf{Set}(\mathcal C) = \textbf{Set}^{\mathcal O_X^\text{op}}$．
预层作为反变函子的作用：

$$
% https://tikzcd.yichuanshen.de/#N4Igdg9gJgpgziAXAbVABwnAlgFyxMJZABgBpiBdUkANwEMAbAVxiRAB12BbOnACwDGjAAQB5APoANAHqccMAB45gENAF8Qa0uky58hFABZyVWoxZs5inACMAZsADKMHBq07seAkTIBGU-TMrIggAKqa2iAYnvo+pADMAebBIABqER563kYJSUFsAGLp7lG6XgbIxv7UgRYhBeEl0VkVvqQATHl1IBmlMdnI8R1dKZqmMFAA5vBEoHYAThBcSGQgOBBIbWb59b0LS0jt1OtIQ9vdnHBMNnAuMACOINQMdDYwDAAKZbEh81iTfBwe0Wy0QxjWG0QAFYaslLOx5EpgPN4Bpnq93l9+gYQH8AUCSvtQUcIUgoYSQadjpDDBSDogAGzUpAAdmeWDAKSgEGuDFY1D4MDoUCQYCYDAYxzoWAYbEgnLGaiAA
\begin{tikzcd}
\mathcal O_X^\text{op} \arrow[rrrr, "F"] &                           &  &    & \textbf{Set}                 \\
U \arrow[dd, "\subseteq"'] \arrow[rrrr]  &                           &  &    & FU                           \\
                                         & {} \arrow[rr, Rightarrow] &  & {} &                              \\
V \arrow[rrrr]                           &                           &  &    & FV \arrow[uu, "\text{res}"']
\end{tikzcd}
$$

其中$FU,FV$是\textbf{截面(section)}，
$\text{res}: FV \to FU$是\textbf{限制(restriction)}．

用$\textbf{cat}$表示由小范畴和函子构成的范畴，
在小范畴中态射集都是集合．
从拓扑空间获取开集范畴的构造如下：

$$
% https://tikzcd.yichuanshen.de/#N4Igdg9gJgpgziAXAbVABwnAlgFyxMJZABgBpiBdUkANwEMAbAVxiRAB12cYAPHAIwBmwACoQ0AXwB6nbn2DiJICaXSZc+QigAs5KrUYs2s3gOEBjOjiUq12PASJkAjPvrNWiEAA1lqkBj2mk6kAMxuhp4gAJp+dhqOOmERHsbsALZWABaWDAAEAPIA+rG2AeoOWsi6rtTuRl6cmTg5jIVFvmWBCVXOpABMKQ0gceVBicihA0NRyvowUADm8ESgggBOEOlIZCA4EEh9BqmNGdm5haMbW0j91PtIU8fDgiDUDHT8MAwAChXBXnWWEWWRwV0220Qdz2B0QAFYytdIU8HohtIiIUgAGz3WEAdneWDAUSgECY-AYrGoWRgdCgSDATAYDHudCwDDYkGJ4Ju8NxSF0IC+YHpiAAtKFdvUooIigAqN4gD5fX7-RIgIEgsEY3mC1Fw6jC0USqWRNLNVr5YqCPIAXjygikCven2+f3GWg1wNBcwkQA
\begin{tikzcd}
\textbf{Top}^\text{op} \arrow[rrrr, "\mathcal O"] &                           &  &    & \textbf{cat}                                               \\
X \arrow[dd, "f"'] \arrow[rrrr]                   &                           &  &    & \mathcal O_X \arrow[dd, "f_*"', bend right]                \\
                                                  & {} \arrow[rr, Rightarrow] &  & {} &                                                            \\
Y \arrow[rrrr]                                    &                           &  &    & \mathcal O_Y \arrow[uu, "\mathcal O_f = f^*"', bend right]
\end{tikzcd}
$$

在连续函数的定义中，
若$f:X \to Y$是连续函数，
那么对于 $Y$ 中的任意开集，
要求其$f$的原像是$X$中的开集，
本质上这就是反变函子的要求，
函子对连续函数$f \ in \textbf{Top}(X,Y)$的作用得到的是拉回
$\mathcal O_f = \text{preim} \ f = f^{-1}$．

借助反变Hom函子，
约定一个靶对象$A$，
可以把对底空间$X$的研究转换为对空间上函数集合$h^A X = \textbf{Set}(X,A)$的研究．
通常希望靶对象$A$具有某种代数结构，
如作为Bool代数的$\textbf{2} = \{0,1\}$，
作为Abel群或交换环的$\mathbb Z$，
或者更强大的，
作为封闭交换域的$\mathbb C$，
这样可以方便地引入代数几何和算子代数的成果．

我们的底空间是拓扑空间$X$，
自然希望限定在连续函数中讨论，
故我们专注于拓扑空间$X$上的复值连续函数代数$C(X) = h^\mathbb C X = \textbf{Top}(X,\mathcal C)$．
在几何上关注拓扑空间范畴$\textbf{Top}^\text{op}$以及特定拓扑空间$X$的开集范畴$\mathcal O_X$，
在代数上关注复交换结合代数$\textbf{CAlg}$和对应的层范畴$\textbf{Sh(CAlg)}$．

从反变Hom函子角度看，
连续函数$f$得到交换代数同态$h^\mathbb C f$：

$$
% https://tikzcd.yichuanshen.de/#N4Igdg9gJgpgziAXAbVABwnAlgFyxMJZABgBpiBdUkANwEMAbAVxiRAB12cYAPHAIwBmwACoQ0AXwB6nbn2DiJICaXSZc+QigAs5KrUYs2s3gOEBhAIIMA5kpVrseAkTIBGffWatEIABrKqiAYTpqupADMnoY+IACagY4aLjqR0d5sABYy7AC2dDiZ-PwABOYlcSUAvGUAFHEAlInB6s5ayLoe1F5Gvtmc+YXFZSV+1XV+TQ4toSnIbqQATOm9IM0hye0RSyuxyvowUDbwRKCCAE4QuUhkIDgQSAsGGX05g0Wl5s0XV0iL1PckNtnqtBCBqAw6PwYAwAAqtMK+c5YGyZHDfS7XRD-O4PRAAVmmPyxwMBiG0RMxSAAbAC8QB2CFYMCxKAQJj8BisaiZGB0KBIMBMBgMAF0LAMNiQFkY37kulIfHUaFgAWIW49WK1fp5AofEaCBq1AC0DXGxpKnAAxlhzlaSmCIVCYfDZloQMjUeiJBQJEA
\begin{tikzcd}
\textbf{Top}^\text{op} \arrow[rrrr, "h^\mathbb C"] &                           &  &    & \textbf{CAlg}                                                      \\
X \arrow[dd, "f"'] \arrow[rrrr]                    &                           &  &    & h^\mathbb C X = C(X)                                               \\
                                                   & {} \arrow[rr, Rightarrow] &  & {} &                                                                    \\
Y \arrow[rrrr]                                     &                           &  &    & h^\mathbb C Y = C(Y) \arrow[uu, "(h^\mathbb C f)(-) = - \circ f"']
\end{tikzcd}
$$

$C(X) = h^\mathbb C X$包含的是全局信息．
若要研究函数集在开集$U \in \mathcal O_X$上的性质，
定义$\mathcal O_U = \mathcal O_X(U) = h^\mathbb C U = \textbf{Top}(U,\mathbb C)$．
包含关系$U \stackrel{\subseteq}{\longrightarrow} V$给出的限制$\mathcal O_V \stackrel{\text{res}}{\longrightarrow} \mathcal O_U$，
使其构成预层，
它满足粘合公理从而构成层．



\subsection{光滑函数层}

第一个应用是微分几何．
拓扑空间$X$赋予无穷阶微分结构构成光滑流形，
相应的我们把连续函数集$C(X)$加强为实值光滑函数集$C^\infty(X,\mathbb R)$，
从而得到微分几何讨论的\textbf{结构层(structure sheaf)}：

$$
% https://tikzcd.yichuanshen.de/#N4Igdg9gJgpgziAXAbVABwnAlgFyxMJZABgBpiBdUkANwEMAbAVxiRAB12BbOnACwDGjAAQB5APoANAHqccMAB45gENAF8Qa0uky58hFABZyVWoxZs5inACMAZsADCAQQYBzDVp3Y8BImQBGU3pmVkQQAFVNbRAMH31-UgBmYPMwkAA1aO89PyNk1NC2ADEM4QBeYUdZdiwwOxwATwAKDIBKbNjdXwNkYyDqEItw4oiKqpq6hpaIjq8u+LzkANIAJkLhkE643N6ktY30zVMYKDd4IlA7ACcILiQyEBwIJBWzIpHOm7ukVepnpD7d6bThwJg2OAweQARxA1AYdBsMAYAAVuglwtcsG4+Dgvrd7ohjE8XogAKyDNKWdjyJTAa7wDTwxHItGLAwgLE4vHzb6Ev4kpBk3kEwH-UmGEU-RAANnFSAA7PC6ukoBBwQxWNQ+DA6FAkGAmAwGP86FgGGxIGBWGoKGogA
\begin{tikzcd}
\mathcal O_X^\text{op} \arrow[rrrr, "F"] &                           &  &    & \textbf{CAlg}                              \\
U \arrow[dd, "\subseteq"'] \arrow[rrrr]  &                           &  &    & FU = C^\infty(U)                           \\
                                         & {} \arrow[rr, Rightarrow] &  & {} &                                            \\
V \arrow[rrrr]                           &                           &  &    & FV = C^\infty(V) \arrow[uu, "\text{res}"']
\end{tikzcd}
$$

在光滑流形的点$x\in X$上做微分是一种局部操作．
不同的全局函数$f,g \in C(X)$可以在$x \in X$的局部相等从而具有相同的微分，
表述为存在$x$的开邻域$U_p \in \mathcal O_X$，
限制在其上$f|_{U_p} = g|_{U_p}$，
构成一个等价关系$f \sim_x g$．
商集$C(X)/\sim_x$为$x$点上的\textbf{函数芽(germ)}．

从层的角度看，
光滑流形$X$得到光滑函数层$\mathcal O_X$，
开集的包含关系转为函数截面的限制关系，
一直限制是一个极限的过程，
直到两个函数在截面上完全相等，
从而具有相同的微分．
点$x\in X$上的\textbf{茎(stalk)}记为$\mathcal O_x = \mathcal O_{X,x} = \lim_{U \ni x} \mathcal O_U$，
它是一个余极限，
得到的$\mathcal O_x$是一个局部化的环，
其中的元素就是函数芽．
作为局部环，
它有唯一的极大理想，
即由$x$处退化的函数芽构成：

$$
\mathfrak m_x = \{ f_x \in \mathcal O_x \ | \ f(x) = 0 \}
$$

其商环

$$
\mathcal O_x/\mathfrak m_x \simeq \mathbb R
$$

切空间有如下同构：

$$
T_x X \simeq \left( \mathfrak m_p / \mathfrak m_p^2 \right)^* = \text{Vct}(\mathfrak m_p / \mathfrak m_p^2,\mathbb R)
$$

这里切向量是极大理想$\mathfrak m_x$上的导子，
$\mathfrak m_p / \mathfrak m_p^2$是余切空间．
余切向量$\text d_x f \in \mathfrak m_p / \mathfrak m_p^2$由$f - f(x)$的等价类给出．


需要视$\mathcal O_X$为\textbf{结构层(structure sheaf)}．


$$
\begin{aligned}
\mathcal O: \textbf{Top}^\text{op} &\to \textbf{Sh(CRing)} \\
X &\mapsto \mathcal O_X
\end{aligned}
$$


$$
\begin{tikzcd}
f \in \mathbf{Sets}(Y,Z) \arrow[rrrr, "h^X"] &  &  &  & \mathbf{Sets}(X,f) \in \mathbf{Sets}(h^X(Y),h^X(Z)) \\
                                              &  & {} \arrow[rr, Rightarrow] &  & {} \\
f \in \mathbf{Sets}(Y,Z) \arrow[rrrr, "F"]   &  &  &  & F(f) \in \mathbf{Sets}(F(Y),F(Z))
\end{tikzcd}
$$

以拓扑空间$X$中的开集为对象，
以开集的包含关系为态射，
构成范畴$\mathcal C = \mathcal O_X$，
其预层为$\textbf{Set}^{\mathcal O_X^\text{op}}$．

令$X$为拓扑空间，
也常用\textbf{结构层(structure sheaf)}的记号来表示拓扑结构$\mathcal O_X$，
视其为$X$上满足拓扑公理的开基族．

考虑在交换域特别是$\mathbb C$上取值的函数全体$h_\mathbb C X = \textbf{Top}(X,\mathbb C)$，
它构成了交换环．
在几何上限制在拓扑空间范畴及连续函数，
在代数上限制在交换环范畴，
预层给出$\mathcal O_X^\text{op} \to \textbf{CRing} :: U \mapsto \mathcal O_U$，
这里的$\mathcal O_U$为预层在开集上的\textbf{截面(section)}．
作为反变函子，
开集的包含嵌入$V \stackrel{\subseteq}{\longrightarrow} U$对应了截面的限制$\mathcal O_U \to \mathcal O_V$．

预层描述了开集的包含关系，进一步要考虑拓扑空间的覆盖和粘合，在预层上发展出层的概念．
这里有一个均衡子正合序列的条件，
对于开集$\forall U\in \mathcal O_X$上的开覆盖$\mathcal U = \{U_i\}$，
一方面若限制对于$\forall U_i \in \mathcal U$有$f|_U = 0$则$f|_U=0$，
即在覆盖上的退化到零对象要求在开集上相应退化；
令一方面若在非空交集$U_i \cap U_j \neq \varnothing$上相等$f_i|_{U_i \cap U_j} = f_j|_{U_i \cap U_j}$，
则$\exists ! f\in \mathcal O_U$ s.t. $f|_{U_i} = f_i$，
这种方式要求交集上两边函数相等，
产生了粘合的效果．

在预层中，开集的包含关系对应到代数结构的态射关系，
开邻域缩小的极限过程对应了代数结构的局部化过程．
在层中，开集交集上代数结构的相容，
确保了整个层可以视为开集上的结构粘合的结果．

可微流形中的\textbf{坐标卡(chart)}和\textbf{图集(atlas)}，
就是拓扑空间中带有局部坐标系的开覆盖．
坐标卡的相容性体现在层中．

\subsection{全局对象}

Topos的定义要求有限极限，
这保证了终对象的存在．
子对象分类器$\Omega$是一个配备了从终对象$\textbf{1}$发出箭头的对象

$$
\text{true}: \textbf{1} \to \Omega
$$

在topos中，
点的概念是通过箭头 

$$
\textbf{1} \to X
$$

来定义的，
箭头称为$X$的全局元素．
在topos中很多对象是没有全局元素的．



\subsection{语言}

量子topos理论将物理系统看作逻辑结构．

$\mathscr L$是一个\textbf{局部语言(local language)}，
若包含如下\textbf{符号(symbol)}：
\begin{itemize}
  \item 
  \textbf{单位类型符号(unit type symbol)} $\textbf{1}$
  和\textbf{真值类型符号(truth value type symbol)} $\Omega$
  \item 可以为空的一组
  \textbf{基本类型符号(ground type symbol)} $\textbf{A},\textbf{B},\textbf{C},\cdots$
  \item 可以为空的一组
  \textbf{函数符号(function symbol)} $\textbf{f},\textbf{g},\textbf{h},\cdots$
\end{itemize}

$\mathscr L$的\textbf{类型符号(type symbol)}递归定义为：
\begin{itemize}
  \item 
  $\textbf{1}$和$\Omega$是类型符号
  \item
  基本类型符号是类型符号
  \item

\end{itemize}

采用直观的ZFC集合论，
并引入Grothendieck宇宙公理，
集合是特定宇宙中的元素，
所有集合都是类．
通过元素性质$P(x)$定义的$\{x:P(x)\}$是类．

Topos中，
任何对象$X$的子对象集合$\text{Sub} \ X$都构成Heyting代数
\textbf{命题语言(propositional language)}描述子对象格．
语法成分：
\begin{itemize}
  \item 原子命题：$\text{Sub} \ X$中的子对象
  \item 逻辑联结词：$\wedge,\vee, \to, \neg$
  \item 常量：最大子对象$\top = 1_X$，最小子对象$\bot = 0_X$
\end{itemize}
PL的逻辑运算直接映射为范畴内的构造，
但它不具备量化能力，
缺少$\forall$和$\exists$，
无法直接描述对象内部的元素关系．

\textbf{高阶语言(higher-order language)}（Mitchell-Bénabou语言）
可以像处理集合论一样处理范畴内的对象．
L是类型语言：
\begin{itemize}
  \item \textbf{类型(type)}：Topos中的每个对象都是一种类型
  \item \textbf{项(term)}：对应范畴中的态射．类型$X$的项$\tau$代表了$X$的广义元素
  \item \textbf{公式(formula)}：其值位于子对象分类器$\Omega$中
\end{itemize}
相比PL，
L引入了高阶逻辑的特性：
\begin{itemize}
  \item 量词：$\forall$和$\exists$，对应于幂对象态射的伴随
  \item 可以定义 $x = _X y$，对应于对角态射 $\Delta: X \to X \times X$
  \item 隶属：可以描述 $x \in A$，其中$A$是幂对象$PX$的一个项
\end{itemize}

Mitchell–Bénabou 语言也称为内语言或类型论语言，
是topos范畴中的一种内部语言，
它允许我们像在集合论中处理集合和元素一样，
在topos中谈论对象和态射．
它是连接范畴论、类型论和逻辑的桥梁．

Topos中的对象不一定有全局元素，
即从终对象出发的态射不一定能区分所有对象．
为了在topos内部推理，
需要一种不依赖全局点的语言．
MB语言提供了一种类型论式的语法，
可以像在集合论中语言使用变量、项和公式，
但语义解释在topos中进行，
使用广义元素即任意对象到目标对象的态射，
代替全局点．

\subsection{量子逻辑}

\cite{BirkhoffvonNeumann1936}
Garrett Birkhoff 和 John von Neumann
它试图通过量子逻辑回答：
量子世界的命题（如“电子的自旋是向上的”）遵循什么样的逻辑规则？

人们希望用直觉主义逻辑来解释量子逻辑，
这需要把Hilbert空间$H \in \textbf{Hilb}$上有闭子空间的集合$\text{Sub} \ H$，
转到某种topos的结构内的子对象分类器．



\subsection{子对象分类器}

用反变Hom函子构造幂集的过程
$h^\textbf 2 : \textbf{Set} \to \textbf{Set} :: X \mapsto \textbf{Set}(X,\textbf 2)$，
幂集记为子对象集$\text{Sub} \ X = PX = \{ A \mid A \subseteq X \}$，
二元集记为$\Omega = \textbf 2 = \{ 0, 1\}$：

$$
\text{Sub} \ X =PX = \{ A \mid A \subseteq X \} \simeq \textbf{Set}(X,\Omega) = \textbf{Set}(X,\textbf 2)
$$

扩展到任意范畴$\mathcal C$，
对象$\Omega \in \mathcal C$使得对于$\forall X \in \mathcal C$：

$$
\text{Sub} \ X \simeq \mathcal C(X,\Omega)
$$

$\Omega$被称为\textbf{子对象分类器(subobject classifier)}．
在topos理论中，
子对象分类器作为逻辑真值的集合，
自带Heyting代数结构。


对照集合$X \in \textbf{Set}$上的幂集$\text{Sub} \ X$，
Hilbert空间$H \in \textbf{Hilb}$上有闭子空间的集合$\text{Sub} \ H$．
如果有一个子对象分类器$\Omega = \textbf{Hilb}$，
它必须是一个Hilbert空间，
且满足

$$
\text{Sub} \ H \simeq\textbf{Hilb}(H,\Omega)
$$

它限制为线性映射的态射集，
$\dim \textbf{Hilb}(H,\Omega) = \dim\Omega \cdot \dim H$

\subsection{CCC}

\cite{Baez2011Rosetta}
给出了范畴的层次，
特别是集合范畴和Hilbert空间范畴在其中的位置：

$$
% https://tikzcd.yichuanshen.de/#N4Igdg9gJgpgziAXAbVABwnAlgFyxMJZABgBpiBdUkANwEMAbAVxiRAB12cYAPHAIwBmwAMJ0cAXxATS6TLnyEUZAIxVajFm07c+Q4AFkCYydNkgM2PASJkATOvrNWiDl14DhAISNgTUmTkrRVtSAGZHTRc3XU9gAAksBn4JAAJOLDB09z1hAGVffzMghRtlUgAWSOdtHLixACccIsCLeWslElIAVmqtVx0PfTyYSWzM7Nj9ERmA80tSzrtSB2onfpih4RFC8TmSjqJliLWo2qntn2M94rbgsuRlqtOagbrpguvTVoXDlArwn1ooNcqIIABbNC7b7zdohf6VIHnLZgyFXPw3H5wh4A3ovDYguIAMSgiWSaQyWUJ0whaE+GO+6hgUAA5vAiKBBA0IUgyCAcBAkCpWlyeYgVNQBUg7CLueDpZLBYgwrKxWFFUgKqr5YgAfyld1tUKNYgAGxGs0mgCcFuW+qQAHYLQ6TSpiBb1fbEAAOC3e13C8yinWmk1OoNyx0m30RsX+r2G2M6q2u91JoV8qXi4UUCRAA
\begin{tikzcd}
\textbf{Cat} \arrow[d]                                   &  &                                           &  &                                          \\
\textbf{MonCat} \arrow[d] \arrow[rrd]                    &  &                                           &  &                                          \\
\textbf{BMonCat} \arrow[d] \arrow[rrd]                   &  & \textbf{CMonCat} \arrow[rrd] \arrow[d]    &  &                                          \\
\textbf{Hilb} \in \textbf{SMonCat} \arrow[d] \arrow[rrd] &  & \textbf{CBMonCat} \arrow[rrd] \arrow[d]   &  & \textbf{CompMonCat} \arrow[d]            \\
\textbf{CartCat} \arrow[d]                               &  & \textbf{CSMonCat} \arrow[rrd] \arrow[lld] &  & \textbf{CompBMonCat} \arrow[d]           \\
\textbf{Set} \in \textbf{CCC}                            &  &                                           &  & \textbf{FdHilb} \in \textbf{CompSMonCat}
\end{tikzcd}
$$

$\textbf{Hilb} \in \textbf{SMonCat}$范畴是一个对称幺半范畴，
它缺少构造CCC乃至topos所需的有限极限．


Bernhard Banaschewskia, and Christopher J. Mulvey.
A globalisation of the Gelfand duality theorem.
Annals of Pure and Applied Logic 137 (2006) 62–103.
\cite{BanaschewskiMulvey2006}

这篇论文的主要贡献是讨论了Grothendieck topos中，
Gelfand对偶定理的构造性版本：

$$
\textbf{CC}^* \simeq \textbf{CCRLoc}^{op}
$$

即交换C*-代数和紧完全正则locale的反变范畴同构．


$$
% https://tikzcd.yichuanshen.de/#N4Igdg9gJgpgziAXAbVABwnAlgFyxMJZABgBpiBdUkANwEMAbAVxiRAB12cYAPHAIwBmwAMIiAvgD1OcHHQBOIcaXSZc+QigAs5KrUYs2nbnyGiRAJQAy4pSpAZseAkTIBGPfWatEIAIIABAC8ASIAFMa8OMAAsnQ8AMpoMADG4gGcAX4AlHaqThpEOh7UXoa+kXyx8Ump6Zl+eQ5qzpokpABMngY+IOExucr56i7and3ebDHBGVxR1YnJabOhYQNKejBQAObwRKCC8hAAtkgd1DgQSADMpT1Gc1Vxi3Ug1Ax0-DAMAAothb4GDBBDgmocTkgAKwXK6IHT6Sa+AAW0nYxzoOCR-H4oTeIA+X1+-1G+OBoKGIHBpzhMLOFKpN1piEh4go4iAA
\begin{tikzcd}
\textbf{CC}^\star                                        &  &  &  & \textbf{CCRL}                                         \\
A = C(\text{MaxSpec} \ A) \arrow[rrrr, "\text{MaxSpec}"] &  &  &  & \text{MaxSpec} \ A \arrow[d]                          \\
C(M) \arrow[u]                                           &  &  &  & M = \text{MaxSpec} \ C(M) \arrow[llll, "h^\mathbb C"]
\end{tikzcd}
$$

1940年代的经典Gelfand理论建立了交换C*-代数和紧Hausdorff空间的对偶．
所有的构造都依赖点和ZFC，
排中律和选择公理都成立．
Banaschewski和Mulvey的工作，
在Grothendieck topos不需要排中律和选择公理．
基于紧完全正则locale范畴$\textbf{CCRLoc}$让拓扑摆脱了对点的依赖，
用开集的frame代替了点集．
使得Gelfand对偶不仅在点集宇宙成立，
也在任何topos宇宙，
如层范畴、
群作用范畴内部成立．

Gelfand 对偶告诉我们，交换 C*-代数等价于拓扑空间（经典物理的相空间）．
量子力学的本质是可观测量的算子代数通常是非交换的．
量子力学可以被看作是非交换拓扑或非交换几何．
在量子理论中，我们不再问点在哪里，
而是通过算子代数来定义系统的状态．

Topos 理论是连接逻辑与几何的桥梁
（通过 Curry-Howard-Lambek 同构）：
Topos 是拓扑空间的推广（层的范畴），
每个 Topos 都有其内部语言（一种类型论）．
在 Topos 中，逻辑不再是排中律必然成立的经典逻辑，而是直觉主义逻辑．
类型论（Type Theory）是这种逻辑的符号化表达．

Topos 量子理论是近年来由 Chris Isham 和 Andreas Döring 等人提出的前沿框架．
既然量子力学因为非交换性而没有经典点空间，
我们可以寻找一个特定的 Topos，
使得在这个 Topos 的内部逻辑下，
量子力学看起来像经典力学．

通过对非交换代数的交换子代数进行层化（Bohr 观点），
Gelfand 谱被推广为 Topos 中的对象．
从而在一个Grothendieck topos中，
用量子态作为类型论中的逻辑项(term)．
量子状态被解释为 Mitchell–Bénabou 语言（一种高阶类型论语言）中的封闭项．


\subsection{附录：DoeringIsham2011}

Andreas Döring and Chris J. Isham.
``What is a Thing?'': 
Topos Theory in the Foundations of Physics.
In New Structures for Physics.
Bob Coecke.
Lecture Notes in Physics,
Vol. 813.
pp. 753-937.
Springer 2011.
\cite{DoeringIsham2011}

\begin{center}
\begin{tabular}{|c|c|c|}
\hline
\textbf{概念} & \textbf{经典物理} & \textbf{量子物理（拓扑斯表示）} \\
\hline
拓扑斯 $\tau$ & $\mathbf{Sets}$ & $\mathbf{Sets}^{\mathcal{V}(H)^{\mathrm{op}}}$ \\
\hline
类型符号 $\Sigma$ 的表示（状态对象） & 辛流形 $S$ & 谱预层 $\underline{\Sigma}$ \\
\hline
类型符号 $R$ 的表示（量值对象） & 实数 $\mathbb{R}$ & 预层 $\underline{R}^{\leftrightarrow}$ 或 $\underline{R}^{\succeq}$ \\
\hline
函数符号 $A: \Sigma\to R$ 的表示 & 函数 $A_\sigma: S\to\mathbb{R}$ & 自然变换 $\check{\delta}(\hat{A}):\underline{\Sigma}\to\underline{R}^{\leftrightarrow}$ \\
\hline
命题表示 & $S$ 的子集 & $\underline{\Sigma}$ 的子对象 \\
\hline
逻辑 & 布尔代数 & 海廷代数（直觉主义逻辑） \\
\hline
\textit{真值} & $\{0,1\}$ & $\Gamma\Omega_{\tau}$ 中的全局元素（筛） \\
\hline
\end{tabular}
\end{center}

语言 $L(S)$ 包含以下基本类型符号：
\begin{itemize}
    \item $\Sigma$：状态对象的前语言符号
    \item $R$：量值对象的前语言符号  
    \item $\Omega$：子对象分类器的前语言符号
    \item $1$：终对象符号
    \item $PT$：类型 $T$ 的幂类型
\end{itemize}

每个物理量对应一个函数符号 $A:\Sigma\to R$．该语言采用\textbf{直觉主义逻辑}的相继式演算作为演绎系统．

在拓扑斯 $\tau$ 中：
\begin{itemize}
    \item 子对象分类器 $\Omega_\tau$ 给出真值对象
    \item 任一对象 $X$ 的子对象集合 $\mathrm{Sub}(X)$ 构成\textbf{海廷代数}
    \item 命题表示为 $\underline{\Sigma}$ 的子对象 $J\hookrightarrow\underline{\Sigma}$
    \item 给定真值对象 $T$，命题的\textit{真值}为 $\nu(J\in T)\in\Gamma\Omega_\tau$
\end{itemize}

与经典布尔逻辑的关键区别：排中律不一定成立，即 $J\vee\neg J \neq \underline{\Sigma}$ 一般情况下．


\begin{itemize}
    \item 对象：$B(H)$ 中的交换冯·诺依曼子代数
    \item 态射：包含关系 $V'\subseteq V$
    \item 每个对象代表一个“语境”或“经典快照”
\end{itemize}

\[\underline{\Sigma}(V) = \text{Gel'fand 谱}(\text{乘法线性泛函}\lambda: V\to\mathbb{C})\]
Kochen-Specker定理 $\iff$ $\underline{\Sigma}$ 无全局元素（无微观态）


\textbf{外具身化}（从上近似）：
\[\delta_o(\hat{P})_V := \bigwedge\{\hat{\alpha}\in P(V) \mid \hat{\alpha}\succeq \hat{P}\}\]

\textbf{内具身化}（从下近似）：
\[\delta_i(\hat{P})_V := \bigvee\{\hat{\beta}\in P(V) \mid \hat{\beta}\preceq \hat{P}\}\]

将命题“$A\in\Delta$”映射为 $\underline{\Sigma}$ 的子对象 $\delta_o(\hat{E}[A\in\Delta])$．

每个自伴算子 $\hat{A}$ 对应自然变换：
\[\check{\delta}(\hat{A})_V(\lambda)(V') := \langle \lambda,\delta(\hat{A})_{V'}\rangle\]
该映射 $\hat{A}\mapsto\check{\delta}(\hat{A})$ 是单射．

\textbf{真值对象} $T_{|\psi\rangle}$（$\mathcal{O}$ 的子对象）：
\[T_{|\psi\rangle_V} := \{\hat{\alpha}\in\mathcal{O}_V \mid \langle\psi|\hat{\alpha}|\psi\rangle = 1\}\]

\textbf{拟态} $w_{|\psi\rangle}$（$\underline{\Sigma}$ 的子对象）：
\[w_{|\psi\rangle_V} := \bigwedge\{\hat{\alpha}\in\mathcal{O}_V \mid |\psi\rangle\langle\psi| \preceq \hat{\alpha}\} = \delta_o(|\psi\rangle\langle\psi|)_V\]

关键等价关系：
\[\text{``}\delta(\hat{P})\in T_{|\psi\rangle}\text{''} \quad\Longleftrightarrow\quad \text{``}w_{|\psi\rangle}\subseteq\delta(\hat{P})\text{''}\]

\begin{enumerate}
    \item \textbf{语境性作为结构性特征}：交换子代数范畴 $\mathcal{V}(H)$ 编码了所有可能的“经典视角”，量子命题的真值在这些视角上变化，形成筛（sieve）．

    \item \textbf{逻辑从外延变为内涵}：经典物理中命题真值由微观态 $s\in S$ 决定；量子物理中真值由真值对象 $T_{|\psi\rangle}$ 和拟态 $w_{|\psi\rangle}$ 决定——真值不在 $\{0,1\}$ 中，而在海廷代数 $\Gamma\Omega$ 中．

    \item \textbf{类型作为范畴中的对象}：形式语言 $L(S)$ 的类型符号（$\Sigma,R,\Omega,1,PT$）在拓扑斯中被解释为对象；函数符号 $A:\Sigma\to R$ 被解释为态射．这使得物理量的“类型”与范畴结构内在统一．

    \item \textbf{量子理论“看起来像”经典理论}：在拓扑斯 $\mathbf{Sets}^{\mathcal{V}(H)^{\mathrm{op}}}$ 内部，量子理论呈现出与经典物理在 $\mathbf{Sets}$ 中完全相同的结构——物理量是态射，命题是子对象，真值由（广义的）态给出．差异仅在于底层的拓扑斯不同．
\end{enumerate}

D\"oring和Isham通过以下对应建立了三者的统一：

\[
\boxed{\text{类型理论（语言 }L(S)\text{）} \xrightarrow{\text{表示}} \text{拓扑斯 }\tau_\varphi(S) \xrightarrow{\text{提供}} \text{逻辑（海廷代数）}}
\]

量子理论在拓扑斯 $\mathbf{Sets}^{\mathcal{V}(H)^{\mathrm{op}}}$ 中的实现表明：
\begin{itemize}
    \item 非交换性被编码为语境范畴的结构
    \item Kochen-Specker定理等价于谱预层无全局元素
    \item 直觉主义逻辑自然出现，排中律不成立
    \item “具身化”（Daseinisation）将量子命题“带入存在”
    \item 真值对象和拟态替代了经典微观态的角色
\end{itemize}

这一框架为量子引力等更深层理论提供了全新的数学工具，同时回应了海德格尔的追问：“什么是物？”——物是在多重语境中“部分存在”的东西．


\printbibliography


\section{复分析}

\subsection{代数}

\subsection{张量代数}

我们用张量代数的商代数来构造~\cite{li2014gaodeng}。
线性空间$X\in\textbf{Vct}$对其自身的张量积，
在同构的意义下满足交换结合律，
故可以记为$T^n X = V \otimes \cdots \otimes V$。
需构造\textbf{张量代数}(tensor algebra) 
$TV = \oplus_{n=0}^\infty T^n V$，
让张量以嵌入方成为张量代数的元$\iota: T^n X \to TX$，
从而让张量积$\otimes: T^m X \times T^n X \to T^{m+n} X$
嵌入到$TX$中封闭。

向量$u,v\in X$嵌入到$TX$，
不混淆的情况下也可以记$u,v \in TX$，
它们的对称性是张量代数中的对易子
$u \otimes v - v \otimes u$刻画的，
用这种对易子生成理想
$I = \langle \{ u \otimes v - v \otimes u \} \rangle$，
构造自然投射$\pi: TV \to TV/I$。
计算$\pi(u) \cdot \pi(v) = u \otimes v + I,\ \pi(v) \cdot \pi(u) = v \otimes u + I$，
注意到$u \otimes v - v \otimes u \in I$，
求差$\varphi(u) \cdot \varphi(v) - \varphi(v) \cdot \varphi(u) = (u \otimes v - v \otimes u) + I = 0$。
于是商代数$SX = TX/I$的所有对易子都退化，
它构成$X$的\textbf{对称代数}(symmetric algebra)．
令$\{x_i\}$为对偶空间$X^* = h^k X$的基，
Zariski-Grothendieck对偶：
对称代数$SX^* = TX^*/I$同构于多项式代数$k[x_i]$，
即交换含幺代数范畴和仿射概形范畴反变等价。

类似的方法，
形如$v \otimes v$的元素生成的理想$I$，
做商代数得到\textbf{外代数}（exterior algebra）
$\wedge X = TX/I$。
$I$ 中的元素具有形式
$ a \otimes (v \otimes v) \otimes b $,
其中 $a, b \in TV,\ v \in V$。
限制在$n$-阶张量空间的 $I_n = I \cap T^n V$，
其中的每个元素都是包含相同相邻元素$v \otimes v$的$n$-阶张量。
外代数可以写成直和：

$$
\wedge X = TX/I 
= \bigoplus_{n=0}^\infty \wedge^n X
= k \oplus X \oplus \bigoplus_{n=2}^\infty T^n X/I_n
$$

自然投影给出了外积：
$v_1 \wedge \cdots \wedge v_p = v_1 \otimes \cdots \otimes v_p + I$。
外积对自身的退化的：
$v \wedge v = v \otimes v + I = I$。
有

\begin{equation}
\dim \wedge^p V = \binom{n}{p} = \frac{n!}{p!(n-p)!}.
\label{eq:dim-exterior-alg}  
\end{equation}

$\{v_1, \cdots, v_k\} \subseteq V$是线性无关的当且仅当
$v_1 \wedge v_2 \wedge \cdots \wedge v_k \neq 0$。

\subsection{多重线性映射}

$n$-重$k$-线性映射$g: X^n \to Y$
称为$n$-\textbf{交错线性映射}（alternating map），
如果对于$\forall v_i \in X$有：
$i \neq j,\ v_i = v_j \ \Rightarrow g(v_i) = 0$。
这样的交错线性映射的全体构成$k$-线性空间，
记为
$\text{Alt}^n(X,Y) \simeq \textbf{Vct}(\wedge^n X,Y)$，
且$\wedge^n X$是$n$-交错映射的泛性质。
当$Y = k$时记为
$\text{Alt}^n(X) = \text{Alt}^n(X,k) = (\wedge^n X)^*$，
它是对偶空间，
其中交错线性映射也称为$n$-\textbf{交错型}（alternating form）。

根据\eqref{eq:dim-exterior-alg}，
列向量空间$k_n$上的$n$-交错线性型空间是$1$-维的：
$\dim \text{Alt}^n k_n = 1$。
定理\cite{li2014gaodeng}：
列向量空间$k_n$只有一个$n$-交错线性型$d\in \text{Alt}^n k_n$，
使得$d(I_n) = 1$，
它就是行列式。

\subsection{双线性型}

在诸多多重线性映射中，
最重要的是双线性型，
它涵盖了一对对偶空间的自然配对，
Hilbert空间的内积，
Riesz表示定理，
Riemann几何的度量，
辛几何的辛结构，
复几何的复结构。
双线性型自动地有二阶张量的矩阵表示，
从而可以转换为一阶输入一阶输出的线性态射。


\subsection{复数}

选择$A = \mathbb C$有诸多的考虑，
首先$\mathbb C$有拓扑，
可以定义连续．
其次复数有自然的对合以便得到C*-代数，
这种对合的本质是伴随，
用实矩阵表示就更容易理解：

$$
z = x + \text i y
=\begin{bmatrix}
  x & -y \\
  y & x
\end{bmatrix}
$$

根据 Riesz 表示定理，
$\forall T\in\mathcal B(H)$，
$\exists! T^*\in\mathcal B(H)$
s.t.
$\forall x, y \in H$，
$\langle T x, y \rangle = \langle x, T^* y \rangle$，
称 $T^*$ 为 $T$ 的 \textbf{伴随算子}．
且满足：
\begin{itemize}
    \item $(T^*)^* = T$，
    \item $(S T)^* = T^* S^*$，
    \item $\|T^*\| = \|T\|$且$\|T^*T\| = \|T\|^2$．
\end{itemize}

矩阵代数都是Banach代数．
实矩阵代数的伴随算子是转置$T^* = T^T$，
复矩阵代数的伴随算子是共轭转置$T^* = \overline{T}^T$．
有时直接用dagger表示伴随$T^* = T^\dagger$．

称$T$为正规算子，若
$T T^* = T^* T$．
矩阵代数中，
正规算子既是正规矩阵，
它在某个标准正交基下可以酉对角化为：
$T = U\Lambda U^{-1}$，
其中$U$为酉矩阵，
$\Lambda$为对角矩阵，
对角元是特征值．
最重要的正规算子是自伴算子和酉算子．

$T = T^*$是自伴算子，
满足 $TT^* = T^2 = T^*T$．
实矩阵代数的自伴算子是对称矩阵，
复矩阵代数的自伴算子是共轭对称矩阵即Hermite矩阵．
自伴算子对角化的特征值是实数$\lambda_i \in \mathbb R$．
自伴算子与内积的对称性相容，
几何上体现了缩放而不旋转．

$T^* = T^{-1}$是酉算子，
满足 $TT^* = I = T^*T$．
实矩阵代数的酉算子是正交矩阵，
复矩阵代数的酉算子是酉矩阵．
酉算子对角化的特征值$\lambda_i = \text e^{\theta_i\text i}$在单位圆．
酉算子保持内积即等距正交，
几何上体现了旋转而不缩放．

线性映射可以分解为缩放和旋转，
自伴算子和酉算子分别描述了这两种几何变换，
它们作为正规算子所描述的这种有几何意义的对称性，
通过对角化完全解耦为一维乘法，
体现了谱的意义．

正规矩阵$A$借助酉矩阵$U$的对角化$A = U\Lambda U^{-1}$，
产生了对角矩阵$\Lambda = \text{diag}(\lambda_i)$．
几何上，
自伴算子的实值$\lambda_i$缩放不旋转，
酉算子的单位复数$\lambda_i = \text e^{\theta_i\text i}$旋转不缩放．

\subsection{二次型}

\subsection{Kähler流形}

在偶数维实线性空间$X \simeq \mathbb R_{2n}$上可以定义辛形式
$\omega \in \mathbb R^{X \times X}$，
它是满足反对称性$\omega(x,y) + w(y,x) = 0$的非退化线性二次型。
还可以定义复结构$J \in \text{End} \ X$，
它是满足$J^2 = - 1_V$的线性自同态。
复结构和辛形式是相容的，若满足：
$J$-不变性$\omega(Jx,Jy) = \omega(x,y)$，
正定性：$\forall x \neq 0:\ \omega(x,Jx) > 0$。
相容性确定了另一个重要的线性二次型——Riemann度量：

$$
g(x,y) = w(x,jy)
$$

三者完全耦合：

$$
\omega(x,y) = g(Jx,y),\quad 
Jx = g^{-1}(\omega(x,-))
$$

把线性代数的这些条件推广到流形$X$上，
辛形式对应辛结构，
复结构对应切空间上的近复结构，
并且加上$J$可积的条件，
就可以定义Kähler流形$(X,g,J,\omega)$，
它统一了复几何的全纯性，
辛几何的Hamilton动力学，
以及Riemann几何的曲率。

\section{C*-代数}

\subsection{乘性算子}

域 $k$ 中的数 $\lambda$ 本身就是 $k$-线性变换，
即 $k \in \operatorname{End} k = \textbf{Vct}_k(k,k)$。
这种$k$-线性变换的直和$\bigoplus_i \lambda_i$
给出了代表伸缩变换的矩阵算子$\bigoplus_i \lambda_i \in \operatorname{End} k^n$。
当 $\lambda_i = 1$ 时，
直和构成了单位矩阵 $I_n$。
这样便把域 $k$ 转为了 $n \times n$ 维 $k$-矩阵代数 $A = (k_n^n, \circ, I_n)$。
为复合的方便从右往左写为：

$$
% https://tikzcd.yichuanshen.de/#N4Igdg9gJgpgziAXAbVABwnAlgFyxMJZAJgBoAGAXVJADcBDAGwFcYkQAPAfQEYACADoCIaFnD7diIAL6l0mXPkIpyFanSat2QxvQC2AIyj1eE00JFjBA3YeNdiZqdPUwoAc3hFQAMwBOEHpIqiA4EEg8NAwsbIggOvpGJvwWosziCXYmUjQ49FiM7Hr0aHBhMpTSQA
\begin{tikzcd}
\lambda_1 x_1 \oplus \lambda_2 x_2 &  & x_1 \oplus x_2 \arrow[ll, "\lambda_1 \oplus \lambda_2", maps to]
\end{tikzcd}
$$

用 $k$ 的直和可以构造线性算子，
反过来我们希望把复杂的线性算子进行直和分解，
若能一直分解到 $1$ 维的数乘，
在数学上就算彻底解决。
这就催生了谱理论，
专门考虑算子 $a$ 在何种条件下相当于 $\lambda_a$ 的乘法。
我们先从最简单的矩阵代数开始。

\subsection{特征多项式}

含幺代数可以讨论逆的问题。
记 $k^\times = k \setminus \{0\}$ 为可逆变换也就是可逆元的全体。
用非退化的 $\lambda_i \in k^\times$ 构造直和，
这样的伸缩变换是可逆的。
$k$-矩阵代数也是含幺代数，
这为 Lie 群的矩阵表示搭建了框架。
$k^\times$ 描述了 $1$ 维可逆 $k$-线性变换，
而 \eqref{eq:det-wedge} 则对应到了任意维度可逆 $k$-线性变换的充分必要条件：
$\det(T) \in k^\times$。
在矩阵代数 $A = (k_n^n,\circ,I_n)$ 中，
行列式刻画了可逆性：
$A^\times = \{ a \in k_n^n \mid \det a \neq 0 \}$。

\subsection{迹}

在$k$-线性空间范畴中有自然配对
$X^* \otimes X \to k: (f \otimes x) = f(x)$，
鉴于：

$$
\text{End} \ X \simeq X^* \otimes X
$$



\section{Gelfand理论}

\subsection{含幺和紧性}

拓扑空间$X\in\textbf{Top}$上的复值连续函数代数
$C(X) = h^\mathbb C = \textbf{Top}(X,\mathbb R)$
是含幺交换代数。
为了使其成为交换C*-代数，
它有自然的复共轭作为对合，
只需要通过赋予范数令其成为Banach代数，
最常见的函数范数为\textbf{上确界范数}（supremum norm）：

$$
\|f\|_\infty = \sup_{x\in X} |f(x)|
$$

若$X$是紧空间，
连续函数必有界，
且$\|f \cdot g \|_\infty \leq \|f\|_\infty \cdot \|f\|_\infty$，
于是$C(X)$构成交换Banach代数和交换C*-代数。
紧致是很高的要求，
为了研究非紧致的空间，
人们转而研究在无穷远处消失的连续函数空间，
往往记为$C_0(X)$，
不混淆的时候我们也用$C(X)$。

关于含幺的问题，
在\textbf{局部紧拓扑空间(locally compact topological space)}范畴，
令$X\in\textbf{LCT}$．
连续函数的全体$\mathbb C^X = \textbf{LCT}(X,\mathbb C)$构成交换C*-代数．
其幺元$\textbf 1$就是在$X$上取值恒为$1$的唯一的函数．
令$C_\infty(X) \subseteq \mathbb C^X = \textbf{LCT}(X,\mathbb C)$
为在无穷远处消失的连续函数全体．
$C_\infty(X)$是交换C*-代数．
只有当$X\in\textbf{CTop}$为紧拓扑空间时，
$\textbf 1 \in C_\infty(X)$，
于是代数的含幺性和拓扑的紧性有了对应．



\subsection{Hausdorff空间}

拓扑空间上的连续函数代数$A = C(X)$是含幺交换Banach代数，
取点$x\in X$，
用点上的求值来定义泛函：

$$
\begin{aligned}
  \varphi_x: C(X) &\to \mathbb C \\
  f &\mapsto \varphi_x(f) = f(x)
\end{aligned}
$$

线性性：
$\varphi_x(af+bg) = (af+bg)(x) = a\varphi_x(f) + b\varphi_x(g)$，
乘性：
$\varphi_x(fg) = (fg)(x) = f(x)g(x) = \varphi_x(f)\varphi_x(g)$，
保幺性：
$\varphi_x(1_A) = 1_A(x) = 1$，
故$\varphi_x$是保幺的乘性线性泛函．

$$
\mathfrak m_x = \text{Ker} \ \varphi_x = \{ f \in C(X) \mid \varphi_x(f) = f(x) = 0 \}
$$

它是$C(X)$上的极大理想，
这样让$X$上的每个点都对应到了$C(X)$上的一个极大理想，
$x \mapsto \mathfrak m_x$是单射，
使得在同构的意义下有嵌入：

$$
X \hookrightarrow \text{Max} \ C(X)
$$

反过来，
$\forall \mathfrak m \in \text{Max} \ C(X)$，
是否能单地映射到$X$上的点，
取决于拓扑空间$X$的拓扑性质．
定义极大理想的零点集

$$
\begin{aligned}
  Z: \text{Max} \ C(X) &\to PX \\
  \mathfrak m &\mapsto Z(\mathfrak m) = \{ x \in X \mid \forall f \in \mathfrak m:\ f(x) = 0 \}
\end{aligned}
$$

以抽象指标$i$标记$X = \{x_i\}$．
若零点集为空，
则$\forall x_i$，
$\exists f_i \in \mathfrak m$
s.t.
$f_i(x_i) \neq 0$．
由于连续性，
存在$x_i$的开领域$U_i$使得$f_i$在其上处处非零．
这样的开邻域构成$X$的开覆盖
$
\mathcal U = \{U_i\}, \qquad
X = \cup_i U_i
$．
若希望在开覆盖的某个子覆盖上构造$X$上的函数：

$$
r(x) = \sum_i | f_i(x) |^2 = \sum_i f_i(x) \overline{f_i(x)} \in \mathfrak m
$$

此时需要求和是有限的，
只有当$X$是紧拓扑空间时，
才存在有限子覆盖，
能够逐点用有限求和来定义$r(x) \in C(X)$．
由于以上$\forall x_i$，
$\exists f_i \in \mathfrak m$
s.t.
$f_i(x_i) \neq 0$．
这样的二次求和使得在任何点上$r(x) > 0$，
使得$r$是可逆函数，
于是$r$不属于某个极大理想．
因此，当$X$是紧拓扑空间时，$Z(\mathfrak m) \neq \varnothing$．

零点集中零点$x_0,x_1$的分离性，
和拓扑空间$X$的分离性相关，
对于分离的零点$x_0 \neq x_1$，
若有函数$f\in C(X)$使得$f(x_0) = f(x_1)$，
那么若有极大理想的零点集$Z(\mathfrak m)$包含$x_0$，
则它也包含$x_1$，
导致$Z: \text{Max} \ C(X) \to PX$的映射得到多于一个零点．

若$X$是Hausdorff空间，
根据Urysohn引理，
存在$f\in C(X)$满足$f(x_0) = 0,\ f(x_1) = 1$，
$\mathfrak m$中在$X_0$处退化的函数构成极大理想
$\mathfrak m_0 \subseteq \mathfrak m$，
且
$h \in \mathfrak m_0$，
又因为极大理想不能真包含，
于是
$\mathfrak m = \mathfrak m_0$．
根据假设条件
$x_1 \in Z(\mathfrak m) = Z(\mathfrak m_0)$，
即$f(x_1) = 0$．
这与$f(x_1) = 1$矛盾，
Hausdorff空间的分离性使得$Z(\mathfrak m)$不能存在一个以上的零点，
此时记$Z(\mathfrak m) = \{ x \}$，
那么$\mathfrak m$就是$x$上的退化函数极大理想．

综合以上，
当$X$是紧Hausdorff空间时，
$Z: \text{Max} \ C(X) \to PX$的映射得到单点集，
即存在$\text{Max} \ C(X) \to PX$的单射嵌入，
使得在同构的意义下有嵌入：

$$
\text{Max} \ C(X) \hookrightarrow X
$$

于是我们建立了同构：

$$
X \simeq \text{Max} \ C(X)
$$

这使得$\text{Max} \ C(X)$构成紧Hausdorff空间．


\subsection{Gelfand-Mazur定理}

设$A = (A,e,0)$是复赋范含幺可除代数，
$e$为单位元，
$0$为零元．
给定$a \in A$，
定义\textbf{预解(resolvent)}函数：

$$
r(\lambda) = (\lambda e - a)^{-1}
$$

它只能定义在$r(\lambda) \neq 0$的地方
称为预解集：

$$
\rho(\lambda) = \{ \lambda \in \mathbb C \ | \ \lambda a - e \neq 0 \}
$$

于是复平面被分为两部分，
另一部分就是谱

$$
\sigma(a) = \{ \lambda \in \mathbb C \ | \ \lambda a - e = 0 \}
$$

谱集是非空紧集．
每个$a \in A$都对应唯一一个复数$\lambda_a$ s.t. $a = \lambda_a e$，
于是：

$$
\sigma(a) = \{ \lambda_a \}
$$

代数 $A$ 中的每个元都只是单位元$e$的复数倍，
这样建立了 $A \simeq \mathbb C$ 的同构．

\subsection{Gelfand表示}


对于任何含幺交换代数$A$，
取极大理想$\mathfrak m \in \text{MaxSpec} \ A$构造商代数，
Gelfand-Mazur定理告诉我们，
含幺交换Banach代数$A \in \textbf{UCB}$对其极大理想$\mathfrak m\in\text{MaxSpec} \ A$做商代数有：

$$
A/\mathfrak m \simeq \mathbb C
$$

$a \in A$作为代表元，
代表了商代数中的元，
根据上式它可以直接视为复数
$a + \mathfrak m \in \mathbb C$．
换句话说我们得到一个二元复值函数：

$$
\begin{aligned}
A \times \text{MaxSpec} \ A &\to A/\mathfrak m \simeq \mathbb C \\
(a,\mathfrak m) &\mapsto a + \mathfrak m 
\end{aligned}
$$

固定一个变元$a \in A$得到极大谱上的复值函数：

\begin{equation}
\begin{aligned}
\varphi_a: \text{MaxSpec} \ A &\to \mathbb C \\
\mathfrak m &\mapsto \varphi_a(\mathfrak m) = a + \mathfrak m 
\end{aligned}
\end{equation}
\label{eq:algebra-ele-as-fun-over-maxspec}

构造这个函数的过程是代数同态，
称为$A$的Gelfand表示：

$$
\begin{aligned}
A &\to \mathbb C^{\text{MaxSpec}} \\
a &\mapsto \varphi_a 
\end{aligned}
$$

\subsection{Gelfand拓扑}

(\ref{eq:algebra-ele-as-fun-over-maxspec})中，
我们不再关注极大理想$\mathfrak m$的内部结构而是将其视为点h，
而代数元素$a$则视为极大理想谱上的复值函数．
为了给这种函数赋予连续性，
需要构造拓扑结构．

令$X\in\textbf{CHaus}$为紧Hausdorff空间，
其上的连续函数代数$C(X)$定义了极大值范数，
这个范数确保$C(X)$是Banach代数．
在Banach代数中，
极大理想是闭的且唯一对应一个特征．
那么有同胚关系：
(张恭庆下p19)

$$
\textbf{CHaus}(\text{MaxSpec} \ C(X) \simeq X)
$$

这是Gelfand对偶性的体现，
对于紧Hausdorff空间，
其几何信息完全被蕴涵在其连续函数代数中．


\subsection{$\Delta(A)$的拓扑}

\subsection{\texorpdfstring{$C(X)$}{C(X)}的拓扑}

在交换Banach代数中，
这意味着乘法线性泛函和极大理想是等价概念，
并且构成拓扑空间$\text{Max} \ A$，
这就是Gelfand谱，
其上可以赋予Gelfand 拓扑，
使其成为局部紧Hausdorff空间．
若$A$含幺则为紧Hausdorff空间．

连续函数代数$C(X)$中有三种拓扑，
它们本质上是同一个拓扑在不同对象上的表现．
设 $X$ 是紧 Hausdorff 空间，
$A = C(X)$ 是 $X$ 上复值连续函数构成的含幺交换 $C^*$-代数，
赋一致范数 $\|f\|_\infty = \sup_{x\in X}|f(x)|$．
记 $\Sigma(A) = \text{Max} \ A$ 为 $A$ 的极大理想空间（即 Gelfand 谱），其上的 Gelfand 拓扑定义为弱 $^*$ 拓扑（视 $\varphi \in \Sigma(A)$ 为 $A$ 上的乘法线性泛函）．

\begin{itemize}
  \item $X$的拓扑：
  $X$ 作为紧 Hausdorff 空间给定的拓扑．
  \item $C(X)$的拓扑：
  由范数 $\|f\|_\infty$ 诱导的拓扑（一致收敛拓扑）．
  \item $\text{Max} \ C(X)$的拓扑：
  在 $\text{Max} \ C(X)$ 上，
  由点态收敛 $\varphi_n \to \varphi \iff \forall f\in C(X),\ \varphi_n(f)\to\varphi(f)$ 定义的拓扑．
\end{itemize}

Gelfand 表示定理给出等距同构
\[
\Gamma: C(X) \xrightarrow{\cong} C(\Sigma(C(X))),\qquad f \mapsto \hat f,\quad \hat f(\varphi) = \varphi(f).
\]
同时，映射
\[
\Phi: X \to \Sigma(C(X)),\qquad x \mapsto \varphi_x,\quad \varphi_x(f)=f(x)
\]
是 \textbf{同胚}．因此：
\[
X \ \text{的原始拓扑} \ \cong \ \Sigma(C(X)) \ \text{上的 Gelfand 拓扑}.
\]

\begin{enumerate}
    \item \textbf{点的收敛}：在 $X$ 中 $x_n \to x$ 当且仅当对每个 $f\in C(X)$，$f(x_n) \to f(x)$，即 $\varphi_{x_n} \to \varphi_x$ 在 Gelfand 拓扑下．
    \item \textbf{Gelfand 拓扑与弱$^*$拓扑}：$\text{Max} \ C(X)$ 作为 $C(X)^*$ 的子集，Gelfand 拓扑正是弱$^*$拓扑，满足 $\varphi_n \to \varphi \iff \forall f,\ \varphi_n(f)\to\varphi(f)$．
    \item \textbf{范数拓扑与 Gelfand 拓扑}：$C(X)$ 的范数拓扑（一致收敛）强于点态收敛，但通过 Gelfand 表示，$C(X)$ 等距同构于 $C(\Sigma(C(X)))$，后者上的范数拓扑由 $\Sigma(C(X))$ 的紧 Hausdorff 拓扑通过最大值范数给出，与 Gelfand 拓扑相容：
    \[
    \|f\|_\infty = \sup_{\varphi\in\Sigma(C(X))} |\hat f(\varphi)|.
    \]
\end{enumerate}

\begin{center}
\begin{tabular}{|c|c|c|}
\hline
\textbf{对象} & \textbf{拓扑} & \textbf{关系} \\
\hline
$X$ & 原始拓扑 & 同胚于 $\text{Max} \ C(X)$ 的 Gelfand 拓扑 \\
\hline
$\text{Max} \ C(X)$ & Gelfand 拓扑 & 弱$^*$拓扑，与 $X$ 同胚 \\
\hline
$C(X)$ & 范数拓扑 & 通过 Gelfand 变换等距同构于 $C(\text{Max} \ C(X))$ \\
\hline
\end{tabular}
\end{center}

在紧 Hausdorff 空间 $X$ 上，
$X$ 的原始拓扑、$\text{Max} \  C(X)$ 的 Gelfand 拓扑以及 $C(X)$ 的范数拓扑通过 Gelfand 对偶达到完美统一：
空间的几何（$X$）、代数（$C(X)$）和谱（$\text{Max} \ C(X)$）三者结构相容，互相决定．

\begin{enumerate}
    \item \textbf{视为函数空间}：对每个 $\varphi \in \Delta(A)$，视 $\varphi: A \to \mathbb{C}$ 为乘法线性泛函（同构于极大理想 $A / \ker \varphi$）．
    
    \item \textbf{初始拓扑}：赋予 $\Delta(A)$ 拓扑 $T$，使得对每个 $a \in A$，求值映射
    \[
    \hat{a}: \Delta(A) \to \mathbb{C}, \quad \hat{a}(\varphi) = \varphi(a)
    \]
    连续．这是 $\Delta(A)$ 上使得所有 $\hat{a}$ 连续的最弱拓扑．
    
    \item \textbf{等价描述}：$T$ 是 $\Delta(A)$ 作为 $A$ 的谱空间（赋值于 $\mathbb{C}$ 的乘法线性泛函子集）在 $\mathbb{C}^A$（所有复值函数）中受**逐点收敛**诱导的**相对拓扑**．
    
    \item \textbf{显式基}：对任意有限集 $\{a_1,\dots,a_n\} \subseteq A$ 和 $\varepsilon > 0$，定义
    \[
    U(\varphi_0; a_1,\dots,a_n; \varepsilon) = \left\{ \varphi \in \Delta(A) \;\middle|\; |\varphi(a_i) - \varphi_0(a_i)| < \varepsilon,\ i=1,\dots,n \right\}.
    \]
    所有这些集合构成 $T$ 的一组邻域基．
    
    \item \textbf{紧致性}：可以证明 $\Delta(A)$ 在该拓扑下是紧 Hausdorff 空间．证明概要：
    \begin{itemize}
        \item $\Delta(A)$ 包含于 $\prod_{a\in A} \overline{B}_{\|a\|}(0)$（由 Banach-Alaoglu 定理，它是紧的）．
        \item 在乘性线性泛函的子集中 $\Delta(A)$ 是闭的（由线性与乘法条件的极限封闭性）．
        \item 故 $\Delta(A)$ 是紧空间的闭子集，从而紧．
    \end{itemize}
    
    \item \textbf{Gelfand 表示}：映射 $\Gamma: A \to C(\Delta(A))$，$a \mapsto \hat{a}$ 是 Banach 代数同态，称为 \textbf{Gelfand 变换}．当 $A$ 是交换 $C^*$-代数时，$\Gamma$ 是等距同构．
\end{enumerate}

Gelfand 拓扑是 $\Delta(A)$ 上使所有 Gelfand 变换 $\hat{a}$ 连续的最弱拓扑．它由 $\mathbb{C}^A$ 的逐点收敛拓扑（即乘积拓扑）限制在 $\Delta(A)$ 上得到，且使 $\Delta(A)$ 成为紧 Hausdorff 空间．该构造建立了交换 $C^*$-代数与紧 Hausdorff 空间之间的对偶．

\subsection{paper}


Bernhard Banaschewskia, and Christopher J. Mulvey.
A globalisation of the Gelfand duality theorem.
Annals of Pure and Applied Logic 137 (2006) 62–103.
\cite{BanaschewskiMulvey2006}

这篇论文的主要贡献是讨论了Grothendieck topos中，
Gelfand对偶定理的构造性版本：

$$
\textbf{CC}^* \simeq \textbf{CCRLoc}^{op}
$$

即交换C*-代数和紧完全正则locale的反变范畴同构．


$$
% https://tikzcd.yichuanshen.de/#N4Igdg9gJgpgziAXAbVABwnAlgFyxMJZABgBpiBdUkANwEMAbAVxiRAB12cYAPHAIwBmwAMIiAvgD1OcHHQBOIcaXSZc+QigAs5KrUYs2nbnyGiRAJQAy4pSpAZseAkTIBGPfWatEIAIIABAC8ASIAFMa8OMAAsnQ8AMpoMADG4gGcAX4AlHaqThpEOh7UXoa+kXyx8Ump6Zl+eQ5qzpokpABMngY+IOExucr56i7and3ebDHBGVxR1YnJabOhYQNKejBQAObwRKCC8hAAtkgd1DgQSADMpT1Gc1Vxi3Ug1Ax0-DAMAAothb4GDBBDgmocTkgAKwXK6IHT6Sa+AAW0nYxzoOCR-H4oTeIA+X1+-1G+OBoKGIHBpzhMLOFKpN1piEh4go4iAA
\begin{tikzcd}
\textbf{CC}^\star                                        &  &  &  & \textbf{CCRL}                                         \\
A = C(\text{MaxSpec} \ A) \arrow[rrrr, "\text{MaxSpec}"] &  &  &  & \text{MaxSpec} \ A \arrow[d]                          \\
C(M) \arrow[u]                                           &  &  &  & M = \text{MaxSpec} \ C(M) \arrow[llll, "h^\mathbb C"]
\end{tikzcd}
$$

1940年代的经典Gelfand理论建立了交换C*-代数和紧Hausdorff空间的对偶．
所有的构造都依赖点和ZFC，
排中律和选择公理都成立．
Banaschewski和Mulvey的工作，
在Grothendieck topos不需要排中律和选择公理．
基于紧完全正则locale范畴$\textbf{CCRLoc}$让拓扑摆脱了对点的依赖，
用开集的frame代替了点集．
使得Gelfand对偶不仅在点集宇宙成立，
也在任何topos宇宙，
如层范畴、
群作用范畴内部成立．

Gelfand 对偶告诉我们，交换 C*-代数等价于拓扑空间（经典物理的相空间）．
量子力学的本质是可观测量的算子代数通常是非交换的．
量子力学可以被看作是非交换拓扑或非交换几何．
在量子理论中，我们不再问点在哪里，
而是通过算子代数来定义系统的状态．

Topos 理论是连接逻辑与几何的桥梁
（通过 Curry-Howard-Lambek 同构）：
Topos 是拓扑空间的推广（层的范畴），
每个 Topos 都有其内部语言（一种类型论）．
在 Topos 中，逻辑不再是排中律必然成立的经典逻辑，而是直觉主义逻辑．
类型论（Type Theory）是这种逻辑的符号化表达．

Topos 量子理论是近年来由 Chris Isham 和 Andreas Döring 等人提出的前沿框架．
既然量子力学因为非交换性而没有经典点空间，
我们可以寻找一个特定的 Topos，
使得在这个 Topos 的内部逻辑下，
量子力学看起来像经典力学．

通过对非交换代数的交换子代数进行层化（Bohr 观点），
Gelfand 谱被推广为 Topos 中的对象．
从而在一个Grothendieck topos中，
用量子态作为类型论中的逻辑项(term)．
量子状态被解释为 Mitchell–Bénabou 语言（一种高阶类型论语言）中的封闭项．



\printbibliography



\end{document}